% =============================================================================
% PNV3425 - Projeto de Sistemas Oceânicos - Trabalho Final - Grupo 02
% RELATÓRIO ÚNICO: caso-base e condição de 10% de lastro
%
% Compilar com XeLaTeX ou Tectonic, dentro da pasta relatorio/:
%     tectonic relatorio.tex
% As figuras são lidas das pastas de resultados das duas versões do projeto.
% Arquivos gerados por scripts (não editar à mão):
%     tabelas_lastro10/*.tex  <- condicao_10pct_lastro/codigo/analises_condicao_10pct_lastro.py
%     trechos/*               <- relatorio/gerar_trechos.py (trechos do código real)
% =============================================================================
\documentclass[11pt,a4paper]{article}

% ---- Fontes: Pagella (texto), Heros (títulos), DejaVu Sans Mono (código) -----
\usepackage{amsmath}
\usepackage{newpxmath}
\usepackage{fontspec}
\setmainfont{texgyrepagella}[Extension=.otf,UprightFont=*-regular,BoldFont=*-bold,
  ItalicFont=*-italic,BoldItalicFont=*-bolditalic]
\setsansfont{texgyreheros}[Extension=.otf,UprightFont=*-regular,BoldFont=*-bold,
  ItalicFont=*-italic,BoldItalicFont=*-bolditalic]
\setmonofont{DejaVuSansMono.ttf}[Scale=MatchLowercase,BoldFont=DejaVuSansMono-Bold.ttf]
\usepackage[brazil]{babel}
\usepackage{microtype}
\usepackage{icomma}

% ---- Página --------------------------------------------------------------------
\usepackage[a4paper,top=2.8cm,bottom=2.6cm,left=2.8cm,right=2.4cm,headheight=14pt]{geometry}
\linespread{1.08}
\setlength{\parindent}{1.2em}
\setlength{\parskip}{0.35em}
\setlength{\emergencystretch}{2em}
\clubpenalty=10000 \widowpenalty=10000 \displaywidowpenalty=10000

% ---- Pacotes gerais --------------------------------------------------------------
\usepackage{graphicx}
\usepackage{booktabs}
\usepackage{array}
\usepackage{multirow}
\usepackage{longtable}
\usepackage{float}
\usepackage{caption}
\usepackage{xcolor}
\usepackage{tikz}
\usetikzlibrary{arrows.meta,positioning}
\usepackage{enumitem}
\usepackage{needspace}
\usepackage{fvextra}
\usepackage{etoolbox}
\usepackage{placeins}
\usepackage{titlesec}
\usepackage{fancyhdr}
\usepackage[most]{tcolorbox}

% ---- Cores da apresentação ---------------------------------------------------
\definecolor{azul}{HTML}{0B2E40}
\definecolor{azulmedio}{HTML}{2E75B6}
\definecolor{laranja}{HTML}{E36C1E}
\definecolor{cinza}{HTML}{5A6B78}

\usepackage[colorlinks=true,linkcolor=azul,citecolor=laranja!80!black,urlcolor=azulmedio,
  bookmarksnumbered=true]{hyperref}

% ---- Títulos: cada seção começa em página nova -------------------------------
\titleformat{\section}{\sffamily\LARGE\bfseries\color{azul}}{\thesection}{0.7em}{}
  [{\vspace{2pt}\color{laranja}\titlerule[1.2pt]}]
\titleformat{\subsection}{\sffamily\large\bfseries\color{azul}}{\thesubsection}{0.7em}{}
\titleformat{\subsubsection}{\sffamily\normalsize\bfseries\color{azul}}{\thesubsubsection}{0.7em}{}
\titleformat{\paragraph}[runin]{\sffamily\normalsize\bfseries\color{azul}}{}{0pt}{}[.]
\titlespacing*{\section}{0pt}{0pt}{1.4em}
\titlespacing*{\subsection}{0pt}{1.6em}{0.6em}
\titlespacing*{\subsubsection}{0pt}{1.2em}{0.4em}
\titlespacing*{\paragraph}{0pt}{0.8em}{0.6em}
\pretocmd{\section}{\clearpage}{}{}
\pretocmd{\subsection}{\FloatBarrier}{}{}

% ---- Cabeçalho e rodapé ----------------------------------------------------------
\renewcommand{\sectionmark}[1]{\markboth{\thesection\quad #1}{}}
\fancypagestyle{principal}{%
  \fancyhf{}%
  \fancyhead[L]{\small\sffamily\color{cinza}PNV3425 · Grupo 02}%
  \fancyhead[R]{\small\sffamily\color{cinza}\nouppercase{\leftmark}}%
  \fancyfoot[C]{\small\sffamily\color{cinza}\thepage}%
  \renewcommand{\headrulewidth}{0.4pt}}
\fancypagestyle{plain}{%
  \fancyhf{}%
  \fancyfoot[C]{\small\sffamily\color{cinza}\thepage}%
  \renewcommand{\headrulewidth}{0pt}}

% ---- Legendas, listas e tabelas -------------------------------------------------
\captionsetup{font=small,labelfont={bf,color=azul},labelsep=period,margin=0.5cm}
\captionsetup[table]{skip=6pt}
\setlist{itemsep=3pt,topsep=5pt,leftmargin=1.6em}
\renewcommand{\arraystretch}{1.18}
\setlength{\abovetopsep}{2pt}
\urlstyle{tt}
\makeatletter\g@addto@macro\UrlBreaks{\do\_}\makeatother

% ---- Caixas ---------------------------------------------------------------------------
% nota de implementação, no fim de cada etapa do modelo
\newtcolorbox{nocodigo}{enhanced,breakable,frame hidden,colback=azul!4,
  borderline west={2.5pt}{0pt}{laranja},arc=0pt,left=8pt,right=8pt,top=4pt,bottom=4pt,
  before skip=10pt,after skip=10pt,fontupper=\small}
\newcommand{\notacodigo}[1]{\begin{nocodigo}\textbf{\sffamily\color{azul}No código.} #1\end{nocodigo}}
% síntese dos resultados de uma seção
\newtcolorbox{sintese}[1][Em resumo]{enhanced,breakable,frame hidden,colback=azulmedio!7,
  borderline west={2.5pt}{0pt}{azul},arc=0pt,left=8pt,right=8pt,top=5pt,bottom=5pt,
  before skip=12pt,after skip=10pt,
  title={#1},coltitle=azul,fonttitle=\sffamily\bfseries,attach title to upper={\quad}}
\newcommand{\lst}[1]{Listagem~\ref{lst:#1}, p.~\pageref{lst:#1}}

% ---- Pastas de figuras (os dois casos têm arquivos com os mesmos nomes) -----
\newcommand{\RB}{../caso_base/resultados}                         % caso-base
\newcommand{\RBC}{../caso_base/resultados/complementar}           % caso-base, análises complementares
\newcommand{\RBA}{../caso_base/resultados/figuras_apresentacao}   % figuras da apresentação
\newcommand{\RL}{../condicao_10pct_lastro/resultados_lastro10}    % condição de 10% de lastro
\newcommand{\RLC}{../condicao_10pct_lastro/resultados_lastro10/complementar}

% ---- Símbolos ---------------------------------------------------------------
\newcommand{\phial}{\varphi_{\mathrm{al}}}
\newcommand{\Hz}{H_{1/3,z}}
\newcommand{\amin}{a_{\text{mín}}}
\newcommand{\Wc}{W_{casco}}

% ---- Tabelas geradas pelo Python: \input primitivo, aceito dentro de tabular -
\makeatletter\newcommand{\linhas}[1]{\@@input #1 }\makeatother

% ---- Listagens de código (trechos copiados do código real) ------------------
\DeclareCaptionType{codigo}[Listagem][Lista de listagens]
\captionsetup[codigo]{justification=raggedright,singlelinecheck=false,margin=0pt}
\fvset{fontsize=\footnotesize,numbers=left,numbersep=6pt,breaklines=true,
       frame=leftline,framerule=0.8pt,rulecolor=\color{laranja!70},xleftmargin=2.4em,
       tabsize=4,numberblanklines=true}
\renewcommand{\theFancyVerbLine}{\sffamily\tiny\color{cinza}\arabic{FancyVerbLine}}
\newcommand{\trechodef}[4]{%
  \expandafter\def\csname trArq#1\endcsname{#2}%
  \expandafter\def\csname trIni#1\endcsname{#3}%
  \expandafter\def\csname trFim#1\endcsname{#4}}
% Gerado por gerar_trechos.py. Não editar à mão.
\trechodef{dados}{caso\_base/codigo/pnv3425\_grupo02.py}{75}{128}
\trechodef{cenario}{caso\_base/codigo/pnv3425\_grupo02.py}{131}{176}
\trechodef{hidrostatica}{caso\_base/codigo/pnv3425\_grupo02.py}{181}{203}
\trechodef{pesos}{caso\_base/codigo/pnv3425\_grupo02.py}{208}{247}
\trechodef{caixas}{caso\_base/codigo/pnv3425\_grupo02.py}{262}{292}
\trechodef{equilibrio}{caso\_base/codigo/pnv3425\_grupo02.py}{295}{354}
\trechodef{alagamento}{caso\_base/codigo/pnv3425\_grupo02.py}{357}{379}
\trechodef{vento}{caso\_base/codigo/pnv3425\_grupo02.py}{384}{420}
\trechodef{criterio}{caso\_base/codigo/pnv3425\_grupo02.py}{423}{475}
\trechodef{heave}{caso\_base/codigo/pnv3425\_grupo02.py}{480}{541}
\trechodef{avaliar}{caso\_base/codigo/pnv3425\_grupo02.py}{546}{598}
\trechodef{lhs}{caso\_base/codigo/pnv3425\_grupo02.py}{665}{690}
\trechodef{problema}{caso\_base/codigo/pnv3425\_grupo02.py}{693}{749}
\trechodef{somapond}{caso\_base/codigo/pnv3425\_grupo02.py}{762}{811}
\trechodef{epsilon}{caso\_base/codigo/pnv3425\_grupo02.py}{814}{857}
\trechodef{fronteira}{caso\_base/codigo/pnv3425\_grupo02.py}{1087}{1102}

\newcommand{\listagem}[2]{\par\needspace{10\baselineskip}\medskip
  \captionof{codigo}[#2]{#2. Arquivo \texttt{\csname trArq#1\endcsname}, linhas
    \csname trIni#1\endcsname\ a \csname trFim#1\endcsname.}\label{lst:#1}%
  \VerbatimInput[firstnumber=\csname trIni#1\endcsname]{trechos/#1.py}\par\medskip}

\newcommand{\integrantes}{Felipe Perpétuo Simões -- 11263255\\[3pt]
  Gustavo Guilherme Silva -- 11820555\\[3pt]
  João Marcos Magalhães Moura -- 12554693}

\begin{document}

% =============================================================================
\hypersetup{pageanchor=false}
\begin{titlepage}
\centering
{\sffamily\large\color{azul} ESCOLA POLITÉCNICA DA UNIVERSIDADE DE SÃO PAULO\par}
{\sffamily\color{cinza} Departamento de Engenharia Naval e Oceânica\par}
\vspace{3cm}
{\sffamily\Large PNV3425 -- Projeto de Sistemas Oceânicos\par}
{\sffamily\color{cinza} Trabalho Final -- 2º semestre de 2026\par}
\vspace{1.4cm}
{\color{laranja}\rule{0.55\textwidth}{1.2pt}\par}
\vspace{0.6cm}
{\sffamily\LARGE\bfseries\color{azul} Projeto Paramétrico de uma Plataforma\\[6pt]
Semissubmersível de Produção\par}
\vspace{0.6cm}
{\color{laranja}\rule{0.55\textwidth}{1.2pt}\par}
\vspace{0.9cm}
{\Large Bloco KZ-2026-B -- Bacia do Baixo Congo (Angola)\par}
\vspace{0.4cm}
{\large Otimização multiobjetivo: peso estrutural do casco $\times$ altura significativa de
\emph{heave}\par}
\vspace{0.3cm}
{\color{cinza} Caso-base e condição de 10\% de lastro\par}
\vfill
{\sffamily\large\bfseries\color{azul} Grupo 02\par}
\vspace{0.4cm}
{\integrantes\par}
\vspace{2cm}
{\sffamily\color{cinza} São Paulo, novembro de 2026\par}
\end{titlepage}
\hypersetup{pageanchor=true}
\pagenumbering{roman}
\pagestyle{plain}

% =============================================================================
\section*{Resumo}
\addcontentsline{toc}{section}{Resumo}

Este relatório documenta o projeto paramétrico de uma plataforma semissubmersível de produção
para o bloco KZ-2026-B, situado em 1.450~m de lâmina d'água na bacia do Baixo Congo, em Angola.
A plataforma tem dois \emph{pontoons} retangulares e quatro colunas quadradas em planta quadrada
de lado $L_p$, e sua planta de processo é dimensionada para o pico de produção de óleo do campo,
de 158.894,48~bpd.

Para cada combinação das seis variáveis de projeto $(L_p, b_p, h_p, b_c, h_c, T)$, um modelo em
Python calcula a hidrostática, os pesos e centros, a curva GZ por projeções 2D, o critério de
áreas para o momento de vento (DNV-OS-C301 e Código MODU 2009), a resposta em \emph{heave} no mar
de 100 anos ($H_s = 5{,}6$~m, $T_p = 15{,}2$~s) e o \emph{air gap} mínimo. O modelo reproduz os
resultados das aulas na plataforma dos Exercícios 7 a 10 com erros entre 0,0\% e 3,7\% e o
deslocamento de três semissubmersíveis reais com erro de até 10\%. A busca pelas melhores
soluções combina 20.000 amostras de hipercubo latino, soma ponderada e método
$\varepsilon$-restrito, ambos resolvidos com SLSQP, e um filtro de não dominância.

No caso-base, a fronteira de Pareto tem 25 pontos, de 10.346~t ($\Hz = 2{,}37$~m) a 23.634~t
($\Hz = 1{,}04$~m). A solução escolhida fica no joelho da fronteira, ponto a partir do qual cada
milímetro de \emph{heave} passa a custar três vezes mais aço: tem $L_p = 111{,}8$~m,
$T = 24{,}0$~m, casco de 14.998~t, $\Hz = 1{,}55$~m, $T_n = 21{,}8$~s e $A/B = 1{,}303$. Em todas
as soluções ótimas o projeto é governado pela estabilidade, com $GM_T = 4$~m e $A/B = 1{,}30$
ativos.

Repetida com lastro mínimo obrigatório de 10\% do volume dos \emph{pontoons}, a otimização leva
a uma solução com 16.026~t de casco (+6,9\%) e 3.611~t de lastro. O KG cai 4,44~m, o
\emph{heave} e o período natural praticamente não mudam, e o deslocamento específico, de 285~t
por mil bopd, chega ao nível da Troll C, a semissubmersível mais enxuta do pôster de 2011.

\vspace{0.8em}
\noindent\textbf{\sffamily\color{azul}Palavras-chave:} semissubmersível; projeto paramétrico;
otimização multiobjetivo; fronteira de Pareto; estabilidade intacta; \emph{heave}; lastro.

\clearpage
\tableofcontents
\clearpage
\listoffigures
\clearpage
\listoftables
\clearpage
\listofcodigos

% =============================================================================
\section*{Lista de símbolos}
\addcontentsline{toc}{section}{Lista de símbolos}

{\small
\begin{longtable}{>{$}l<{$}p{11.6cm}}
\toprule
\text{Símbolo} & Significado [unidade] \\
\midrule
\endhead
L_p & lado da planta quadrada, comprimento dos \emph{pontoons} e lado do convés [m] \\
b_p,\ h_p & largura e altura dos \emph{pontoons} [m] \\
b_c,\ h_c & lado e altura das colunas quadradas [m] \\
T & calado [m] \\
a & \emph{air gap}: distância da linha d'água ao fundo do convés, $a = h_p + h_c - T$ [m] \\
h_d,\ h_m & altura do \emph{deck box} e altura média dos módulos de \emph{topside} [m] \\
d_p & distância entre os eixos dos \emph{pontoons}, $d_p = L_p - b_p$ [m] \\
l_d & comprimento livre dos \emph{pontoons} entre colunas, $l_d = L_p - 2b_c$ [m] \\
\nabla,\ \Delta & volume e massa deslocados [m$^3$, t] \\
KB,\ KG & alturas do centro de carena e do centro de gravidade sobre a quilha [m] \\
BM_T,\ BM_L & raios metacêntricos transversal e longitudinal [m] \\
GM_T,\ GM_L & alturas metacêntricas transversal e longitudinal [m] \\
A_{wp} & área de linha d'água, $4b_c^2$ [m$^2$] \\
\varphi & ângulo de inclinação [$^\circ$ ou rad] \\
GZ(\varphi),\ HA(\varphi) & braço de endireitamento e braço de emborcamento do vento [m] \\
\phial & ângulo de alagamento (fundo do convés toca a água na borda) [$^\circ$] \\
\varphi_1,\ \varphi_2 & primeiro e segundo interceptos entre GZ e HA [$^\circ$] \\
A,\ B & áreas sob GZ e sob HA de $0$ a $\phial$ [m$\cdot$rad] \\
z_{CLR} & altura do centro de resistência lateral sobre a quilha [m] \\
HM_0,\ HM_{ud} & momentos de vento das parcelas laterais e da face inferior do convés [kN$\cdot$m] \\
C_s,\ C_h & coeficientes de forma e de altura do Código MODU 2009 \\
\omega,\ k & frequência da onda [rad/s] e número de onda em águas profundas, $k = \omega^2/g$ [1/m] \\
A_{33},\ B_{33},\ C_{33} & massa adicional [t], amortecimento [kN$\cdot$s/m] e rigidez [kN/m] em \emph{heave} \\
C_A & coeficiente 2D de massa adicional do retângulo (DNV-RP-C205) \\
\omega_n,\ T_n & frequência [rad/s] e período natural [s] de \emph{heave} \\
F_3/\zeta_a & força de excitação vertical por amplitude de onda [kN/m] \\
RAO & operador de amplitude de resposta em \emph{heave} [m/m] \\
S_\zeta,\ S_z & espectros do mar e da resposta em \emph{heave} [m$^2\cdot$s] \\
m_{0z},\ \Hz & momento de ordem zero [m$^2$] e altura significativa de \emph{heave} [m] \\
H_s,\ T_p,\ \gamma & altura significativa [m], período de pico [s] e fator de pico do JONSWAP \\
\xi & fração do amortecimento crítico em \emph{heave} \\
\eta_{max},\ z_{min},\ \amin & crista máxima, posição mais baixa do convés e \emph{air gap} mínimo [m] \\
\Wc & peso estrutural do casco (\emph{pontoons} + colunas), primeiro objetivo [t] \\
Q,\ \text{BOE} & produção de pico [mil boe/dia] e [boe/dia] \\
w_1,\ \varepsilon & peso do \emph{heave} na soma ponderada e teto de $\Hz$ no $\varepsilon$-restrito \\
V_{lastro},\ V_{pontoons} & volume de lastro e volume interno dos \emph{pontoons}, $2L_pb_ph_p$ [m$^3$] \\
\bottomrule
\end{longtable}}
\addtocounter{table}{-1}

% =============================================================================
\clearpage
\pagenumbering{arabic}
\pagestyle{principal}
\section{Introdução}\label{sec:intro}

A Kwanza Offshore Energy adquiriu, por 22 anos, o direito de exploração do bloco KZ-2026-B, em
águas profundas da bacia do Baixo Congo, ao largo de Angola. Este trabalho dimensiona a
plataforma semissubmersível de produção de óleo do campo pelo método paramétrico desenvolvido
nos Exercícios 7 a 10 da disciplina e escolhe, de forma justificada, uma solução final. O método
segue as referências do enunciado~\cite{chakrabarti,goncalves,brotzman,gallala,dsouza,poster}.

O projeto persegue dois objetivos ao mesmo tempo: o menor peso estrutural do casco e o menor
movimento significativo de \emph{heave}. Os dois são conflitantes, porque reduzir o \emph{heave}
exige afundar os \emph{pontoons}, o que pede calado maior e colunas mais altas e, portanto, um
casco mais pesado. Não existe, assim, uma única plataforma ótima, e sim um conjunto de soluções
de compromisso, a fronteira de Pareto no plano peso $\times$ $\Hz$. A solução final é escolhida
sobre essa fronteira.

O repositório do grupo guarda duas versões do projeto, e este relatório trata das duas. O
\textbf{caso-base} (pasta \path{caso_base/}) segue a letra do enunciado, que exige apenas lastro
positivo, e ocupa as Seções~\ref{sec:dados} a~\ref{sec:sensibilidade}. A \textbf{condição de 10\%
de lastro} (pasta \path{condicao_10pct_lastro/}) repete o projeto inteiro com a restrição
adicional $V_{lastro} \ge 0{,}10\,V_{pontoons}$ e é apresentada na Seção~\ref{sec:lastro}. Como o
modelo, os dados, os limites e os objetivos são os mesmos, o código da segunda versão difere do
primeiro apenas nas linhas marcadas com \texttt{[LASTRO 10\%]} (\lst{diff}).

A ordem do relatório acompanha a da apresentação final do grupo, como mostra a
Tabela~\ref{tab:mapa}. Em cada etapa, o texto parte da teoria, segue para a aplicação ao campo,
com o procedimento e os resultados, e termina com uma nota que indica o trecho de código que faz
o cálculo. Esses trechos são cópias exatas do código do repositório, com a numeração de linhas
original, e estão reunidos no Apêndice~\ref{ap:codigo}. Por fim, a Seção~\ref{sec:discussao}
discute as hipóteses de entrada que o enunciado não fixa (BOE com ou sem gás, peso da planta,
limites das variáveis e alturas do convés) e as limitações do modelo.

\begin{table}[!htb]
\centering
\caption{Correspondência entre a apresentação final e o relatório.}\label{tab:mapa}
\small
\begin{tabular}{rll}
\toprule
Slide & Tema & Seção \\
\midrule
2 & O campo KZ-2026-B & \ref{sec:dados} \\
3 & Modelo paramétrico: geometria e variáveis & \ref{sec:modelo} \\
4 & Validação: plataforma dos Exercícios 7 a 10 & \ref{sec:validacao} \\
5 & Procedimento de variação sistemática & \ref{sec:procedimento} \\
6 e 7 & Conjunto viável, fronteira de Pareto e vista lateral & \ref{sec:pareto} \\
8 & Por que a solução intermediária & \ref{sec:escolha} \\
9 & Ficha técnica comparativa dos três cenários & \ref{sec:cenarios} \\
10 & Comparativo com plataformas reais & \ref{sec:reais} \\
11 a 13 & Solução final: pesos, estabilidade, \emph{heave} e \emph{air gap} & \ref{sec:final} \\
14 & Sensibilidade às hipóteses & \ref{sec:sensibilidade} \\
15 & Impacto do lastro mínimo operacional & \ref{sec:lastro} \\
16 & Conclusões & \ref{sec:conclusoes} \\
\bottomrule
\end{tabular}
\end{table}

% =============================================================================
\section{O campo KZ-2026-B e os dados de projeto}\label{sec:dados}

Antes de dimensionar a plataforma, é preciso traduzir os dados do campo em requisitos de
projeto: a capacidade da planta, o mar e o vento de projeto e as correlações que ligam a
produção ao peso e à área do convés. Esta seção reúne esses dados e as hipóteses adotadas onde o
enunciado não é explícito.

\subsection{Campo e reservatório}

O bloco fica a 190~km da costa, em 1.450~m de lâmina d'água e solo argiloso
(Tabela~\ref{tab:campo}). O reservatório tem 2,2 bilhões de barris de óleo original, dos quais
660 milhões são recuperáveis com o fator de recuperação de 30\%, além de 30 bilhões de m$^3$ de
gás recuperável.

\begin{table}[!htb]
\centering
\caption{Características do bloco KZ-2026-B.}\label{tab:campo}
\begin{tabular}{ll}
\toprule
Item & Valor \\
\midrule
Operador e concessão & Kwanza Offshore Energy, 22 anos \\
Localização & Bacia do Baixo Congo, Angola \\
Distância até a costa & 190 km \\
Lâmina d'água & 1.450 m \\
Solo marinho & argiloso \\
Volume original de óleo & 2,2 bilhões de barris, 31 °API \\
Fator de recuperação & 30\% (660 milhões de barris recuperáveis) \\
Gás recuperável & 30 bilhões de m$^3$ \\
Reservatório & topo a 2.700 m abaixo do leito, coluna de 120 m \\
\bottomrule
\end{tabular}
\end{table}

\subsection{Curva de produção e capacidade da planta}

A capacidade de uma unidade de produção é definida pelo pico da curva de produção, porque a
planta de processo precisa tratar a maior vazão diária ao longo da vida do campo. No KZ-2026-B,
esse pico é de 158.894,48~bpd e ocorre entre os anos 5 e 7 (Figura~\ref{fig:producao}). A
integral da curva dá 659 milhões de barris, coerente com a reserva recuperável de 660 milhões, o
que confirma a leitura da curva. A planta é, portanto, dimensionada para
\begin{equation}
\text{BOE} = 158.894{,}48~\text{boe/dia}, \qquad Q = \frac{\text{BOE}}{1000} = 158{,}894~\text{mil boe/dia}.
\end{equation}

\begin{figure}[!htb]
\centering
\includegraphics[width=0.66\textwidth]{\RBA/f_producao.png}
\caption[Curva de produção de óleo do campo]{Curva de produção de óleo do campo KZ-2026-B, com o
pico dos anos 5 a 7.}\label{fig:producao}
\end{figure}

O enunciado fala em plataforma de produção de óleo e fornece a curva de produção de óleo, e por
isso o caso-base usa apenas essa produção. Se o gás também fosse processado a bordo, os 30
bilhões de m$^3$ equivaleriam a cerca de 183 milhões de boe (com 5,8 mil pés$^3$ por boe, ou
164,2~m$^3$/boe), 27,7\% a mais que o óleo recuperável, e o pico passaria a 202,9 mil boe/dia se o
gás seguisse o mesmo perfil de produção. Esse caso é tratado como sensibilidade na
Seção~\ref{sec:sensibilidade}, e a escolha é discutida na Seção~\ref{sec:disc-boe}.

\subsection{Condições ambientais}

A Tabela~\ref{tab:ambiente} resume as condições extremas do local. O projeto usa o mar de 100
anos, com $H_s = 5{,}6$~m e $T_p = 15{,}2$~s.

\begin{table}[!htb]
\centering
\caption{Condições ambientais extremas omnidirecionais do campo.}\label{tab:ambiente}
\begin{tabular}{lccc}
\toprule
Período de retorno & $H_s$ [m] & $T_p$ [s] & Vento [m/s] \\
\midrule
1 ano   & 3,4 & 11,5 & 12,0 \\
10 anos & 4,5 & 13,4 & 18,0 \\
100 anos & \textbf{5,6} & \textbf{15,2} & \textbf{24,0} \\
\bottomrule
\end{tabular}
\end{table}

Para descrever esse mar, a DNV-RP-C205~\cite{c205} recomenda escolher o fator de pico do espectro
JONSWAP pela razão $T_p/\sqrt{H_s}$: $\gamma = 5$ até 3,6, $\gamma = \exp(5{,}75 - 1{,}15\,
T_p/\sqrt{H_s})$ entre 3,6 e 5 e $\gamma = 1$ (Pierson-Moskowitz) acima de 5. Como no campo
$T_p/\sqrt{H_s} = 6{,}42$, a recomendação levaria a $\gamma = 1$. O projeto adota o valor usado na
disciplina, $\gamma = 3{,}3$, e informa também o resultado com $\gamma = 1$; na solução final, o
primeiro dá o maior $\Hz$ (1,552 contra 1,454~m).

O vento entra no projeto apenas pelo critério de estabilidade, que usa as velocidades do Código
MODU~\cite{modu}: 51,5~m/s para tempestade severa (100 nós) e 36~m/s para operação normal
(70 nós). Ambas ficam bem acima do vento de 100 anos do local, de 24~m/s. Como o momento de vento
cresce com $V^2$, o vento local produziria só $(24/51{,}5)^2 = 22\%$ do momento de tempestade, de
modo que a verificação é conservadora para o campo. Já a corrente, de 0,95~m/s na superfície, não
entra no projeto paramétrico, que não inclui a ancoragem nem o \emph{offset} da plataforma.

\subsection{Dados fixos e correlações da disciplina}\label{sec:fixos}

Três correlações ligam a produção às dimensões e aos pesos da plataforma. A primeira dá a área de
convés necessária,
\begin{equation}
A_{conv,\text{mín}} = 42{,}246\,Q + 4850{,}1 = 11.562{,}8~\text{m}^2
\quad\Rightarrow\quad L_p \ge \sqrt{11.562{,}8} = 107{,}53~\text{m},
\end{equation}
que é o único limite do domínio tirado diretamente dos dados do campo
(Seção~\ref{sec:dominio}). A segunda dá o peso da planta de produção, 0,083~t por boe/dia de
capacidade (83~kg/boe), ou seja, 13.188~t. A estrutura do convés e a carga operacional valem,
cada uma, 35\% desse peso (4.616~t), e o \emph{topside} soma 22.420~t. O valor de 83~kg/boe vem
do código original do grupo; a Seção~\ref{sec:disc-planta} mostra que ele reproduz o KG da
plataforma de referência das aulas e a mediana de \emph{topside} das semissubmersíveis reais. A
terceira correlação atribui aos \emph{risers} e à amarração 4\% e 3\% do deslocamento, aplicados
no fundo do convés.

Faltam ainda as alturas do \emph{deck box} ($h_d$) e dos módulos ($h_m$), que o enunciado não
fornece e que foram estimadas por semelhantes. No pôster de 2011~\cite{poster}, a Independence Hub
tem convés de 140 × 220 × 36~pés, isto é, 11~m de altura para dois níveis, e as unidades listadas
têm de 1 a 3 níveis de convés. O valor adotado, $h_d = 8$~m, fica abaixo dessa referência; já
$h_m = 12$~m não consta do pôster e é uma estimativa do grupo. Como os módulos respondem por cerca
de 54\% do momento de vento das parcelas laterais na solução final, as duas alturas são objeto de
análise de sensibilidade (Seção~\ref{sec:sens-hdhm}). A Tabela~\ref{tab:fixos} reúne todos os
dados fixos.

\begin{table}[!htb]
\centering
\caption{Dados fixos do projeto.}\label{tab:fixos}
\begin{tabular}{llr}
\toprule
Dado & Símbolo & Valor \\
\midrule
Produção de pico (só óleo) & BOE & 158.894,48 boe/dia \\
Área mínima de convés & $42{,}246Q + 4850{,}1$ & 11.562,8 m$^2$ \\
Planta de produção & 0,083 t/boe & 13.188 t \\
\emph{Topside} (planta + 35\% + 35\%) & $W_{top}$ & 22.420 t \\
Massa específica da água do mar & $\rho$ & 1,025 t/m$^3$ \\
Massa específica do ar & $\rho_{ar}$ & 1,222 kg/m$^3$ \\
Peso específico estrutural & $\gamma_e$ & 0,270 t/m$^3$ \\
\emph{Risers} / amarração & -- & 4\% / 3\% de $\Delta$ \\
Altura do \emph{deck box} / dos módulos & $h_d$ / $h_m$ & 8 / 12 m (estimadas) \\
Vento de tempestade / operação & $V$ & 51,5 / 36 m/s \\
Mar de projeto & $H_s$, $T_p$ & 5,6 m, 15,2 s \\
Amortecimento em \emph{heave} & $\xi$ & 0,05 \\
Fator de pico do JONSWAP & $\gamma$ & 3,3 \\
\bottomrule
\end{tabular}
\end{table}

\notacodigo{os dados fixos, os coeficientes do Código MODU e as hipóteses parametrizadas estão na
\lst{dados}; a classe \texttt{Cenario}, com a área mínima e a crista máxima, e o domínio
\texttt{BOUNDS} estão na \lst{cenario}.}


% =============================================================================
\section{Modelo paramétrico}\label{sec:modelo}

O modelo paramétrico transforma seis dimensões principais em todas as grandezas de que o projeto
precisa para verificar as restrições e calcular os dois objetivos. Esta seção descreve a
geometria, o domínio das variáveis e, em seguida, cada etapa do cálculo, na ordem em que o código
as executa: hidrostática, pesos, estabilidade em grandes ângulos, vento, \emph{heave}, espectro de
resposta e \emph{air gap}. Ao final, reúne as restrições, as correções feitas no modelo inicial do
grupo e as hipóteses adotadas.

\subsection{Geometria e variáveis de projeto}

A plataforma tem dois \emph{pontoons} retangulares que correm na direção longitudinal $x$, com
comprimento $L_p$, e quatro colunas quadradas nos cantos, centradas sobre os \emph{pontoons}. O
convés é quadrado, com o mesmo lado $L_p$. As seis variáveis de projeto,
\begin{equation}
x = (L_p,\ b_p,\ h_p,\ b_c,\ h_c,\ T),
\end{equation}
aparecem na Figura~\ref{fig:geometria}. Delas decorrem a distância entre os eixos dos
\emph{pontoons}, $d_p = L_p - b_p$, as coordenadas dos centros das colunas,
$x_c = \pm(L_p - b_c)/2$ e $y_c = \pm(L_p - b_p)/2$, o calado nas colunas, $t_c = T - h_p$, a área
de convés, $L_p^2$, e o \emph{air gap}, $a = h_p + h_c - T$, que é a distância da linha d'água ao
fundo do convés. A geometria só é consistente se a linha d'água cortar as colunas; por isso o
modelo impõe também $T - h_p \ge 1$~m, restrição que não fica ativa em nenhuma solução ótima.

\begin{figure}[!htb]
\centering
\includegraphics[width=0.85\textwidth]{\RBA/f_geometria_param.png}
\caption[Geometria paramétrica da plataforma]{Geometria paramétrica: planta e vista lateral, com
as seis variáveis de projeto, o \emph{air gap} $a$ e as alturas do \emph{deck box} e dos módulos.
LA é a linha d'água de projeto.}\label{fig:geometria}
\end{figure}

\subsection{Domínio das variáveis}\label{sec:dominio}

A Tabela~\ref{tab:limites} mostra os limites adotados ao lado da faixa das semissubmersíveis do
pôster de 2011~\cite{poster}, apresentada na Seção~\ref{sec:reais}. O limite inferior de $L_p$ é
a raiz quadrada da área mínima de convés. Os demais vêm do código original do grupo, já que o
enunciado não fixa limites para as variáveis. Em todas as variáveis o domínio é mais amplo que a
frota real e, portanto, não restringe a busca de forma artificial. O preço disso é que os
extremos da fronteira de Pareto podem encostar nos limites e corresponder a cascos sem precedente
na frota, efeito discutido na Seção~\ref{sec:disc-limites}.

\begin{table}[!htb]
\centering
\caption[Domínio das variáveis de projeto]{Domínio das variáveis de projeto e faixa das
semissubmersíveis do pôster de 2011.}\label{tab:limites}
\begin{tabular}{lrr}
\toprule
Variável & Limite adotado & Faixa no pôster (mín.--máx.) \\
\midrule
$L_p$ [m] & 107,53--150 & -- \\
$b_p$ [m] & 8--25 & 9,1--19,8 \\
$h_p$ [m] & 6--20 & 6,1--10,1 \\
$b_c$ [m] & 8--25 & 12,2--19,8 \\
$h_c$ [m] & 20--50 & 33,2--43,9 \\
$T$ [m] & 12--40 & 20,4--32,0 \\
\bottomrule
\end{tabular}
\end{table}

\subsection{Fluxo de cálculo}

Para cada $x$, o modelo percorre as etapas da Figura~\ref{fig:fluxo}, e cada uma usa as saídas da
anterior. Uma avaliação completa leva cerca de 6~ms, dos quais só 0,7~ms correspondem às etapas
que não dependem da curva GZ. Essa velocidade é o que torna viável a busca sistemática com dezenas
de milhares de avaliações descrita na Seção~\ref{sec:procedimento}.

\begin{figure}[!htb]
\centering
\begin{tikzpicture}[node distance=4mm and 6mm,
  box/.style={draw=azul!60,rounded corners,fill=azul!5,minimum width=5.4cm,minimum height=7mm,
              align=center,font=\small},
  arr/.style={-Stealth,thick,azul!70}]
  \node[box,fill=laranja!15] (x) {Variáveis $x = (L_p, b_p, h_p, b_c, h_c, T)$};
  \node[box,below=of x] (h) {1. Hidrostática: $\nabla$, $\Delta$, KB, BM, KM};
  \node[box,below=of h] (p) {2. Pesos e centros: $\Wc$, $W_{lastro}$, KG, GM};
  \node[box,below=of p] (c) {3. Convés e \emph{air gap}: $L_p^2$, $a$};
  \node[box,below=of c] (g) {4. Curva GZ 2D (0--50$^\circ$), $\phial$};
  \node[box,below=of g] (v) {5. Vento: $HM_0$, $HM_{ud}$, $HA$, $A/B$};
  \node[box,right=12mm of h] (f) {6. \emph{Heave}: $F_3$, $A_{33}$, $C_{33}$, $T_n$, RAO};
  \node[box,below=of f] (s) {7. Espectro: $S_z$, $m_{0z}$, $\Hz$};
  \node[box,below=of s] (ag) {8. \emph{Air gap} mínimo: $\eta_{max}$, $z_{min}$, $\amin$};
  \node[box,below=of ag,fill=red!8] (r) {Restrições e objetivos\\$\Wc$, $\Hz$};
  \node[box,below=of r,fill=green!10] (o) {Otimização: LHS + soma ponderada\\+ $\varepsilon$-restrito $\to$ Pareto};
  \draw[arr] (x)--(h); \draw[arr] (h)--(p); \draw[arr] (p)--(c); \draw[arr] (c)--(g);
  \draw[arr] (g)--(v); \draw[arr] (h.east)--(f.west); \draw[arr] (f)--(s); \draw[arr] (s)--(ag);
  \draw[arr] (ag)--(r); \draw[arr] (v.east)--(r.west); \draw[arr] (r)--(o);
\end{tikzpicture}
\caption{Fluxo de cálculo do modelo paramétrico.}\label{fig:fluxo}
\end{figure}

\notacodigo{a função \texttt{avaliar()} encadeia as etapas e \texttt{tabela\_restricoes()} monta as
restrições (\lst{avaliar}).}

% -----------------------------------------------------------------------------
\subsection{Hidrostática (Exercício 7)}\label{sec:hidro}

Pelo princípio de Arquimedes, a plataforma em equilíbrio desloca uma massa de água igual ao seu
peso, e o centro de carena B é o centroide do volume submerso. Em pequenas inclinações, B se
desloca sobre um arco de raio $BM = I/\nabla$, em que $I$ é o momento de inércia da área de linha
d'água em relação ao eixo de inclinação. O metacentro M fica, então, a $KM = KB + BM$ da quilha, e
a estabilidade inicial é medida pela altura metacêntrica $GM = KM - KG$.

Na semissubmersível, os \emph{pontoons} ficam inteiramente submersos e as colunas, submersas até
o calado $t_c = T - h_p$. O volume deslocado e a altura do centro de carena são
\begin{gather}
\nabla = 2 L_p b_p h_p + 4 b_c^2 t_c, \qquad \Delta = \rho \nabla, \\
KB = \frac{2 L_p b_p h_p \, (h_p/2) + 4 b_c^2 t_c \, (h_p + t_c/2)}{\nabla} .
\end{gather}
Como a linha d'água corta só as quatro colunas, $A_{wp} = 4b_c^2$, e o teorema dos eixos paralelos
dá
\begin{equation}
I_T = 4\left(\frac{b_c^4}{12} + b_c^2 y_c^2\right), \quad
I_L = 4\left(\frac{b_c^4}{12} + b_c^2 x_c^2\right), \quad
BM_{T,L} = \frac{I_{T,L}}{\nabla}, \quad KM_{T,L} = KB + BM_{T,L} .
\end{equation}
Nessas expressões, o termo $b_c^2 y_c^2$ domina: a estabilidade inicial de uma semissubmersível
vem do afastamento entre as colunas, e não da área de cada uma. Na solução final, por exemplo,
$I_T = 1{,}53\times 10^6$~m$^4$ e $BM_T = 39{,}1$~m.

\notacodigo{função \texttt{hidrostatica()} (\lst{hidrostatica}).}

% -----------------------------------------------------------------------------
\subsection{Pesos, centros e altura metacêntrica (Exercício 8)}\label{sec:pesos}

O peso total da plataforma é a soma dos itens da Tabela~\ref{tab:itenspeso}, e o KG é a média das
alturas dos centros ponderada pelos pesos. O lastro fecha o balanço: é a parcela do deslocamento
que não é ocupada por casco, \emph{topside}, \emph{risers} e amarração.

\begin{table}[!htb]
\centering
\caption{Itens de peso e altura do centro (VCG, a partir da quilha).}\label{tab:itenspeso}
\begin{tabular}{lll}
\toprule
Item & Peso & VCG \\
\midrule
Planta de produção & $0{,}083~\text{t/boe} \times \text{BOE}$ & $h_p + h_c + h_d + h_m/2$ \\
Estrutura do convés & $0{,}35\, W_{planta}$ & $h_p + h_c + h_d/2$ \\
Carga operacional & $0{,}35\, W_{planta}$ & $h_p + h_c + h_d/2$ \\
\emph{Pontoons} & $\gamma_e \, 2 L_p b_p h_p$ & $h_p/2$ \\
Colunas & $\gamma_e \, 4 b_c^2 h_c$ & $h_p + h_c/2$ \\
\emph{Risers} & $0{,}04\, \Delta$ & $h_p + h_c$ \\
Amarração & $0{,}03\, \Delta$ & $h_p + h_c$ \\
Lastro & $\Delta - \sum W$ & $n/2$ \\
\bottomrule
\end{tabular}
\end{table}

O peso estrutural de \emph{pontoons} e colunas é o peso específico estrutural aplicado ao volume
envolvido, e a soma das duas parcelas é o primeiro objetivo do projeto:
\begin{equation}
\Wc = W_{pontoons} + W_{colunas} = \gamma_e\left(2 L_p b_p h_p + 4 b_c^2 h_c\right).
\end{equation}
O lastro é água nos \emph{pontoons}. Com nível de enchimento $n = V_{lastro}/(2 L_p b_p)$, seu
centro fica em $n/2$, e não em $h_p/2$, de modo que pouco lastro fica muito perto da quilha. Esse
detalhe pesa na condição de 10\% (Seção~\ref{sec:lastro}). Com todos os itens,
\begin{equation}
KG = \frac{\sum W_i z_i}{\Delta}, \qquad GM_{T,L} = KM_{T,L} - KG ,
\end{equation}
sem correção de superfície livre. O enunciado exige $1~\text{m} \le GM_T \le 4$~m: o limite
inferior garante estabilidade inicial com margem, e o superior evita uma plataforma rígida demais
em \emph{roll}, com período natural curto e acelerações altas no convés.

\notacodigo{função \texttt{pesos()} (\lst{pesos}).}

% -----------------------------------------------------------------------------
\subsection{Curva GZ pelo método de projeções 2D (Exercício 9)}\label{sec:gz-teoria}

O GM descreve a estabilidade apenas em ângulos pequenos. Em ângulos maiores, a forma do volume
submerso muda: o \emph{pontoon} do lado alto sai da água, e o convés pode tocar a água do lado
baixo. A estabilidade passa a ser descrita pelo braço de endireitamento $GZ(\varphi)$, a distância
horizontal entre o centro de gravidade e o centro de carena na posição inclinada. Em cada ângulo,
a plataforma afunda ou sobe até que o volume submerso volte a ser $\nabla$.

O Exercício 9 calcula essa curva projetando a plataforma em cada plano principal e representando-a
por retângulos com espessura fora do plano. No plano transversal (inclinação em torno de $x$), há
dois \emph{pontoons} de largura $b_p$ e espessura $L_p$ e duas colunas de largura $b_c$ e espessura
$2b_c$, já que as duas colunas de cada bordo ficam uma atrás da outra, centradas em $\pm y_c$. No
plano longitudinal (inclinação em torno de $y$), os dois \emph{pontoons} se sobrepõem em um único
retângulo $L_p \times h_p$ de espessura $2b_p$, e as colunas, de largura $b_c$ e espessura $2b_c$,
ficam nas extremidades. Como todos os elementos são prismas na direção normal ao plano e a linha
d'água inclinada não varia nessa direção, a projeção é exata para esta geometria de caixas; o
volume do \emph{deck box} não é incluído.

No referencial do corpo, com $u$ a coordenada no plano e $z$ a altura sobre a quilha, a linha
d'água inclinada é a reta $z = \alpha + u\tan\varphi$. Para cada ângulo, o parâmetro $\alpha$ é a
raiz de
\begin{equation}
V(\alpha, \varphi) = \sum_i t_i\, A_i^{sub}(\alpha, \varphi) = \nabla ,
\end{equation}
em que $t_i$ é a espessura e $A_i^{sub}$ a área do retângulo $i$ abaixo da reta, calculada de forma
exata, assim como os seus momentos. A raiz é obtida pelo método de Newton com a derivada exata
\begin{equation}
\frac{\partial V}{\partial \alpha} = \sum_i t_i\, \ell_i ,
\end{equation}
em que $\ell_i$ é o comprimento, na direção $u$, do trecho da linha d'água dentro do retângulo
$i$: subir a reta de $d\alpha$ acrescenta uma faixa de área $\ell_i\,d\alpha$. Como $V(\alpha)$ tem
quinas quando a linha d'água passa por um canto, o Newton é protegido por um intervalo de
bissecção (Seção~\ref{sec:cod-desempenho}). Conhecido o centro de carena $(u_B, z_B)$ e com G em
$(0, KG)$,
\begin{equation}
GZ(\varphi) = u_B \cos\varphi + (z_B - KG)\sin\varphi .
\end{equation}
Quando $\varphi \to 0$, a inclinação $dGZ/d\varphi$ tende ao GM da hidrostática, propriedade usada
na validação (Seção~\ref{sec:validacao}).

A curva vale até o ângulo de alagamento $\phial$, em que o fundo do convés na borda da plataforma
($u = L_p/2$, $z = h_p + h_c$) toca a água:
\begin{equation}
h_p + h_c - \left(\alpha(\phial) + \tan\phial\,\frac{L_p}{2}\right) = 0 .
\end{equation}
A equação é resolvida pelo método de Brent, com $\alpha(\varphi)$ recalculado a volume constante.
Depois de $\phial$ a curva perde o significado físico, porque o modelo não inclui o volume do
convés, e por isso o critério de áreas só a usa até esse ângulo.

\notacodigo{os retângulos da projeção e as integrais exatas estão em \texttt{caixas\_2d()} e
\texttt{\_integrais\_submersas()} (\lst{caixas}); o equilíbrio a volume constante e a curva GZ,
em \texttt{equilibrio\_2d()} e \texttt{curva\_gz()} (\lst{equilibrio}); o ângulo de alagamento, em
\texttt{angulo\_alagamento()} (\lst{alagamento}).}

% -----------------------------------------------------------------------------
\subsection{Momento de vento e critério de áreas (Exercício 9)}\label{sec:vento-teoria}

O vento empurra a parte emersa da plataforma, e a água reage sobre a parte submersa. As duas
forças formam um binário cujo braço vai do centro de pressão do vento ao centro de resistência
lateral, o centroide da silhueta submersa vista na direção do vento. O critério de áreas da
DNV-OS-C301~\cite{c301} e do Código MODU 2009~\cite{modu} para unidades estabilizadas por colunas
compara os dois efeitos em termos de energia: a área sob a curva de endireitamento, até o segundo
intercepto ou até o ângulo de alagamento, o que ocorrer antes, deve superar em pelo menos 30\% a
área sob a curva do vento até o mesmo ângulo. A folga de 30\% cobre rajadas e o balanço da
plataforma.

Para montar a curva do vento, o modelo calcula a pressão dinâmica $q = \frac12 \rho_{ar} V^2$ e a
força em cada parcela exposta,
\begin{equation}
F_i = q\, C_s\, C_h\, A_i ,
\end{equation}
com as parcelas da Tabela~\ref{tab:parcelas}. Os coeficientes de forma são $C_s = 1{,}0$ para
colunas, \emph{deck box} e face inferior do convés e $C_s = 1{,}1$ para os módulos agrupados. O
coeficiente de altura $C_h$ vem da tabela do Código MODU 2009 em função da altura do centroide
sobre a linha d'água: 1,00 até 15,3~m, 1,10 até 30,5~m, 1,20 até 46~m e assim por diante. Seguindo
o Exercício 9, as quatro colunas são consideradas expostas, sem efeito de sombra, e os módulos
ocupam toda a largura $L_p$.

\begin{table}[!htb]
\centering
\caption[Parcelas de vento]{Parcelas de vento. As alturas $z_i$ são medidas a partir da linha
d'água.}\label{tab:parcelas}
\begin{tabular}{llll}
\toprule
Parcela & Área exposta & $z_i$ & Momento \\
\midrule
Colunas (4, sem sombra) & $4 b_c a$ & $a/2$ & \multirow{3}{*}{$HM_0$} \\
\emph{Deck box} & $L_p h_d$ & $a + h_d/2$ & \\
Módulos & $L_p h_m$ & $a + h_d + h_m/2$ & \\
Face inferior do convés & $L_p^2$ & $a$ & $HM_{ud}$ \\
\bottomrule
\end{tabular}
\end{table}

O braço de cada parcela vai do seu centroide ao centro de resistência lateral, cuja altura sobre a
quilha depende do plano:
\begin{gather}
z_{CLR,T} = \frac{L_p h_p \,(h_p/2) + 2 b_c t_c\,(h_p + t_c/2)}{L_p h_p + 2 b_c t_c}, \qquad
z_{CLR,L} = \frac{2 b_p h_p \,(h_p/2) + 2 b_c t_c\,(h_p + t_c/2)}{2 b_p h_p + 2 b_c t_c}, \\
b_i = z_i + (T - z_{CLR}), \qquad HM_0 = \sum_{\text{3 parcelas}} F_i b_i, \qquad HM_{ud} = F_{ud}\, b_{ud}.
\end{gather}
No plano transversal, a silhueta lateral dos \emph{pontoons} ($L_p h_p$) puxa o $z_{CLR}$ para
baixo e alonga os braços, e por isso esse plano tende a ser o crítico. Quando a plataforma inclina,
a área lateral projetada diminui com $\cos\varphi$ e a face inferior do convés passa a receber
vento proporcionalmente a $\sin\varphi$, o que dá o braço de emborcamento
\begin{equation}
HA(\varphi) = \frac{HM_0\cos\varphi + HM_{ud}\sin\varphi}{\Delta g} .
\end{equation}

O critério compara, então, as áreas sob as duas curvas, com $\varphi$ em radianos e integração
pela regra dos trapézios:
\begin{equation}
A = \int_0^{\phial} GZ\, d\varphi, \qquad B = \int_0^{\phial} HA\, d\varphi, \qquad A \ge 1{,}3\,B .
\end{equation}
Os interceptos $\varphi_1$, que é o equilíbrio estático sob o vento, e $\varphi_2$ são as raízes
de $GZ - HA$, obtidas por interpolação linear. A plataforma atende ao critério quando, nos quatro
casos (2 planos $\times$ 2 velocidades), vale $A \ge 1{,}3B$ e a GZ é positiva entre $0^\circ$ e
$\varphi_2$, ou entre $0^\circ$ e $\phial$ se $\varphi_2$ não ocorre antes. Essa última
verificação usa o indicador $\min(GZ/\sin\varphi)$, que não se anula artificialmente em
$\varphi = 0$ e vale $GM$ para ângulos pequenos. O código integra sempre até $\phial$; como em
todas as soluções relatadas $\varphi_2 > \phial$, esse é o limite correto, e, se fosse
$\varphi_2 < \phial$, o trecho a mais teria $GZ < HA$ e reduziria $A/B$, o que torna a escolha
conservadora.

\notacodigo{o centro de resistência lateral e as forças, braços e momentos por plano estão em
\texttt{z\_clr()} e \texttt{tabela\_vento()} (\lst{vento}); o critério de áreas e o laço sobre os
dois planos, em \texttt{criterio\_areas()} e \texttt{estabilidade()} (\lst{criterio}).}

% -----------------------------------------------------------------------------
\subsection{\emph{Heave} (Exercício 10)}\label{sec:heave-teoria}

Em ondas regulares de amplitude $\zeta_a$ e frequência $\omega$, o \emph{heave} $z(t)$ segue a
equação linear de um sistema de um grau de liberdade,
\begin{equation}
(\Delta + A_{33})\,\ddot z + B_{33}\,\dot z + C_{33}\,z = F_3\, e^{i\omega t} ,
\end{equation}
com massa adicional $A_{33}$, amortecimento $B_{33}$ e rigidez hidrostática
$C_{33} = \rho g A_{wp}$. A frequência e o período naturais são
\begin{equation}
\omega_n = \sqrt{\frac{C_{33}}{\Delta + A_{33}}}, \qquad T_n = \frac{2\pi}{\omega_n} .
\end{equation}
Como a área de linha d'água de uma semissubmersível é pequena, $T_n$ é longo. O enunciado exige
$T_n > 20$~s, acima do período de pico do mar ($T_p = 15{,}2$~s), para manter a ressonância fora
da faixa de maior energia das ondas.

A massa adicional usa o coeficiente 2D de retângulo da DNV-RP-C205~\cite{c205}, $C_A(b_p/h_p)$,
aplicado a uma seção circular de raio $b_p/2$ e ao comprimento livre dos \emph{pontoons} entre
colunas, $l_d = L_p - 2b_c$. O amortecimento é uma fração $\xi = 0{,}05$ do crítico:
\begin{equation}
A_{33} = 2\, l_d\, C_A\, \rho\, \pi \left(\frac{b_p}{2}\right)^2, \qquad
B_{33} = 2\xi\sqrt{C_{33}(\Delta + A_{33})} .
\end{equation}

A força de excitação tem duas origens. A primeira é a força de Froude-Krylov, a integral da
pressão da onda não perturbada sobre as superfícies horizontais molhadas. Em águas profundas, essa
pressão decai com a profundidade $d$ como $\rho g \zeta_a e^{-kd}$, com $k = \omega^2/g$. Ela atua
para cima no fundo dos \emph{pontoons} (área $A_p = 2L_pb_p$, profundidade $T$) e para baixo no
topo dos \emph{pontoons} fora das colunas (área $A_p - A_{wp}$, profundidade $T - h_p$).
Reagrupando, surgem uma parcela ``das colunas'', $\rho g A_{wp} e^{-k(T-h_p)}$, positiva, e uma
``dos \emph{pontoons}'', $-\rho g A_p (e^{-k(T-h_p)} - e^{-kT})$, negativa. A segunda origem é a
difração, estimada pela massa adicional multiplicada pela aceleração vertical da água no centro
dos \emph{pontoons}, $-\omega^2\zeta_a e^{-k(T-h_p/2)}$ (aproximação de onda longa). Para ondas
que se propagam na direção transversal $y$, os dois bordos ficam em $y = \pm d_p/2$; suas forças
se defasam de $\pm k d_p/2$, e a soma ganha o fator $\cos(k d_p/2)$:
\begin{equation}
\frac{F_3}{\zeta_a} = \cos\!\left(\frac{\omega^2 d_p}{2g}\right)\left[
\rho g A_{wp} e^{-k(T-h_p)}
- \rho g\, A_p \left(e^{-k(T-h_p)} - e^{-kT}\right)
- A_{33}\,\omega^2 e^{-k(T-h_p/2)} \right].
\end{equation}

Essa expressão explica o comportamento típico de uma semissubmersível. Para $\omega \to 0$,
$F_3/\zeta_a \to \rho g A_{wp} = C_{33}$, e a plataforma acompanha a onda longa. À medida que
$\omega$ cresce, a parcela das colunas decai devagar, enquanto as dos \emph{pontoons} crescem com
$\omega^2$. Em uma frequência de cancelamento $\omega_c$, as parcelas se equilibram e $F_3 = 0$;
se $\omega_c$ fica perto de $\omega_n$, a própria força de excitação amortece a ressonância. A
resposta em regime permanente por amplitude de onda é, por fim,
\begin{equation}
RAO(\omega) = \frac{|F_3/\zeta_a|}{\left|C_{33} - (\Delta + A_{33})\omega^2 + i B_{33}\omega\right|} .
\end{equation}

% -----------------------------------------------------------------------------
\subsection{Espectro de resposta (Exercício 10)}\label{sec:espectro-teoria}

O mar irregular é modelado como uma soma de ondas regulares com fases aleatórias, descrita pelo
espectro $S_\zeta(\omega)$. Como o sistema é linear, cada componente é multiplicada pelo RAO, e o
espectro da resposta é $S_z = RAO^2 S_\zeta$. Para um processo de banda estreita, as amplitudes
seguem a distribuição de Rayleigh, e a altura significativa, média do terço maior, vale
$4\sqrt{m_0}$.

No modelo, o espectro do mar é o JONSWAP da DNV-RP-C205~\cite{c205}, com $\omega_p = 2\pi/T_p$ e
$\sigma = 0{,}07$ para $\omega \le \omega_p$ e 0,09 acima:
\begin{equation}
S_\zeta(\omega) = (1 - 0{,}287\ln\gamma)\,\frac{5}{16} H_s^2 \omega_p^4 \omega^{-5}
\exp\!\left[-\frac54\left(\frac{\omega}{\omega_p}\right)^{-4}\right]
\gamma^{\exp\left[-\frac{(\omega - \omega_p)^2}{2\sigma^2\omega_p^2}\right]} ,
\end{equation}
\begin{equation}
S_z(\omega) = RAO^2(\omega)\, S_\zeta(\omega), \qquad m_{0z} = \int_0^\infty S_z\, d\omega,
\qquad \Hz = 4\sqrt{m_{0z}} .
\end{equation}
As integrais são feitas pela regra dos trapézios em 1.500 frequências entre 0,02 e 2,5~rad/s,
malha fina o bastante para resolver o pico de ressonância. No pico do mar do campo,
$\omega_p = 0{,}413$~rad/s, $k = 0{,}0174$~m$^{-1}$ e o comprimento de onda é de 361~m, cerca de
três vezes o lado da plataforma.

\notacodigo{o coeficiente de massa adicional, o espectro JONSWAP e o cálculo completo do \emph{heave}
estão em \texttt{ca\_retangulo()}, \texttt{espectro\_mar()} e \texttt{heave()} (\lst{heave}).}

% -----------------------------------------------------------------------------
\subsection{\emph{Air gap} mínimo}\label{sec:airgap-teoria}

Um mar de 3 horas tem cerca de 1.000 ondas. Para alturas com distribuição de Rayleigh, a maior
altura mais provável em $N$ ondas é $H_{max} \approx H_s\sqrt{\ln N/2}$, o que dá
$H_{max} \approx 1{,}86\,H_s$ para $N \approx 1.000$, e, na teoria linear, a crista é metade da
altura. O mesmo raciocínio vale para o \emph{heave}. O \emph{air gap} mínimo soma a crista máxima
à descida máxima da plataforma:
\begin{equation}
\eta_{max} = \frac{1{,}86\, H_s}{2}, \qquad z_{min} = -\frac{1{,}86\,\Hz}{2}, \qquad
\amin = \eta_{max} - z_{min}, \qquad a \ge \amin .
\end{equation}
Para o mar do campo, $\eta_{max} = 5{,}208$~m. A soma supõe que a crista máxima coincide com a
posição mais baixa da plataforma, o que ignora a fase dada pelo RAO e é conservador.

% -----------------------------------------------------------------------------
\subsection{Restrições}\label{sec:restricoes}

As grandezas das etapas anteriores alimentam as restrições da Tabela~\ref{tab:restricoes}, que
definem o conjunto viável do problema de otimização.

\begin{table}[!htb]
\centering
\caption{Restrições do problema de otimização.}\label{tab:restricoes}
\begin{tabular}{ll}
\toprule
Restrição & Expressão \\
\midrule
Área mínima de convés & $L_p^2 \ge 42{,}246\,Q + 4850{,}1 = 11.562{,}8$ m$^2$ \\
Estabilidade inicial & $1~\text{m} \le GM_T \le 4~\text{m}$ \\
Período natural de \emph{heave} & $T_n > 20$ s \\
Lastro & $W_{lastro} > 0$ e $V_{lastro} \le 0{,}9\, V_{pontoons}$ \\
Colunas e \emph{pontoons} & $b_c \le b_p$ \\
Critério de áreas & $\min(A/B) \ge 1{,}3$ e $GZ > 0$, nos 4 casos \\
\emph{Air gap} & $a \ge \amin$ \\
Consistência geométrica (adicional) & $T - h_p \ge 1$ m \\
\midrule
Só na condição de 10\% de lastro (Seção~\ref{sec:lastro}) & $V_{lastro} \ge 0{,}10\, V_{pontoons}$ \\
\bottomrule
\end{tabular}
\end{table}

No código, cada restrição é escrita como $g_j(x) \ge 0$ e normalizada por uma escala própria.
Durante a otimização, todas recebem uma margem de $10^{-4}$, as estritas recebem mais $10^{-3}$ e o
$A/B$ recebe mais 0,005, porque a GZ é calculada com passo de $1^\circ$ durante a busca. A margem
da restrição de lastro positivo faz o otimizador parar em 11~t de lastro, valor que aparece em
todas as soluções sem lastro. Uma restrição é marcada como \emph{ativa} quando é atendida com folga
menor que 1\% do limite.

% -----------------------------------------------------------------------------
\subsection{Revisão do modelo inicial do grupo}\label{sec:revisao}

O código parte da versão inicial do grupo, revisada antes da otimização para alinhar o modelo aos
Exercícios 9 e 10 e tornar a busca confiável. As mudanças estão marcadas no código com
\texttt{[REV n]} e resumidas na Tabela~\ref{tab:rev}.

\begin{table}[!htb]
\centering
\caption{Itens da revisão do modelo inicial.}\label{tab:rev}
\small
\begin{tabular}{>{\raggedright\arraybackslash}p{2.2cm}>{\raggedright\arraybackslash}p{4.6cm}>{\raggedright\arraybackslash}p{7.4cm}}
\toprule
Item & Versão inicial & Versão final \\
\midrule
{[REV 1]} Vento & área de 2 colunas e braço referido a $T/2$ &
4 colunas sem sombra ($4b_ca$), como no Exercício 9; braço até o $z_{CLR}$ de cada plano;
tabela de vento montada por plano \\
{[REV 2]} \emph{Heave} & sem fator de defasagem; massa adicional com $L_p$ &
fator $\cos(\omega^2 d_p/2g)$ com $d_p = L_p - b_p$; massa adicional com $l_d = L_p - 2b_c$;
$\gamma = 3{,}3$, com $\gamma = 1$ também reportado \\
{[REV 3]} Otimização & falhas do SLSQP voltavam em silêncio ao ponto de partida &
só aceita o resultado do SLSQP se for viável; \texttt{success} registrado; continuação entre
pesos; método $\varepsilon$-restrito para os trechos não convexos \\
{[REV 4]} Hipóteses & fixas no código &
lastro mínimo, gás no BOE e $h_d$/$h_m$ parametrizados por variável de ambiente \\
{[REV 5]} Validação & sem comparação com as aulas & tabela de erros contra os valores das aulas \\
{[REV 6]} Rótulos e $\gamma$ & $\gamma$ fixado na definição da função &
$\gamma$ lido na chamada; rótulos dos gráficos com quebra de linha \\
Desempenho & linha d'água por bissecção; \texttt{pandas} dentro do laço &
Newton com salvaguarda, chute inicial no calado e tabelas só no relatório
(Seção~\ref{sec:cod-desempenho}) \\
\bottomrule
\end{tabular}
\end{table}

\subsection{Hipóteses e limitações do modelo}\label{sec:hipoteses}

Além das hipóteses do enunciado, o modelo adota as seguintes simplificações:
\begin{enumerate}
\item a massa adicional usa o $C_A$ de retângulo da DNV-RP-C205, em função de $b_p/h_p$, e o
comprimento livre $l_d = L_p - 2b_c$, desprezando a contribuição das colunas;
\item a excitação em \emph{heave} considera ondas de través, com o fator $\cos(kd_p/2)$, a força de
Froude-Krylov da onda não perturbada e a difração pela aproximação de onda longa;
\item o espectro é o JONSWAP com $\gamma = 3{,}3$, e o resultado com $\gamma = 1$ é reportado;
\item o vento atua sobre quatro colunas sem sombra, com braços até o $z_{CLR}$ de cada plano,
módulos em toda a largura $L_p$ e coeficientes $C_s$ e $C_h$ do Código MODU 2009;
\item o lastro fica no nível real de enchimento dos \emph{pontoons}, sem superfície livre;
\item a GZ vem da projeção 2D, sem o volume do \emph{deck box} e sem acoplamento com o trim.
\end{enumerate}
O enunciado também limita o cálculo do \emph{air gap}: a coincidência entre crista máxima e posição
mais baixa ignora a fase dada pelo RAO, \emph{roll} e \emph{pitch} são desprezados, e efeitos não
lineares, como cristas acima de $H_{max}/2$ e a amplificação pelas colunas, ficam de fora; eles
seriam tratados pela DNVGL-OTG-13~\cite{otg13}. A Seção~\ref{sec:disc-limitacoes} discute o efeito
dessas simplificações nos resultados.


% =============================================================================
\section{Validação do modelo}\label{sec:validacao}

Antes de usar o modelo para otimizar, é preciso verificar se ele reproduz resultados conhecidos. A
validação foi feita em três frentes: a comparação com os resultados das aulas na plataforma dos
Exercícios 7 a 10, um conjunto de verificações internas de consistência e a aplicação do modelo a
semissubmersíveis reais.

\subsection{Plataforma dos Exercícios 7 a 10}

A plataforma de referência tem $L_p = 100$~m, $b_p = h_p = b_c = 15$~m, $h_c = 45$~m e
$T = 30$~m e foi avaliada com BOE = 120.000, $H_s = 12{,}2$~m, $T_p = 14{,}6$~s e $h_m = 15$~m,
valor usado no Exercício 9. A Tabela~\ref{tab:valref} reúne os resultados do modelo, e a
Tabela~\ref{tab:valaula} os compara com os valores obtidos em aula.

\begin{table}[!htb]
\centering
\caption{Resultados do modelo para a plataforma de referência.}\label{tab:valref}
\begin{tabular}{lr@{\hspace{1.2cm}}lr}
\toprule
Grandeza & Modelo & Grandeza & Modelo \\
\midrule
$\nabla$ [m$^3$] & 58.500 & $GM_T = GM_L$ [m] & 5,825 \\
$\Delta$ [t] & 59.962,5 & $\phial$ transversal / longitudinal [$^\circ$] & 27,58 / 30,40 \\
KB [m] & 10,962 & $z_{CLR}$ transversal / longitudinal [m] & 10,96 / 15,00 \\
$BM_T = BM_L$ [m] & 28,077 & $A_{33}$ [t] & 38.291 \\
$\Wc$ [t] & 23.085 & $\Hz$ com $\gamma = 1$ [m] & 2,983 \\
$W_{planta}$ [t] & 9.960 & $\eta_{max}$ [m] & 11,35 \\
$W_{lastro}$ [t] & 15.748 & $a$ / $\amin$ [m] & 30,0 / 14,12 \\
\bottomrule
\end{tabular}
\end{table}

\begin{table}[!htb]
\centering
\caption{Comparação com os resultados obtidos em aula.}\label{tab:valaula}
\begin{tabular}{lrrr}
\toprule
Grandeza & Modelo & Aula & Erro \\
\midrule
KG [m] & 33,213 & 33,00 & +0,6\% \\
$\min A/B$ transversal & 2,310 & 2,33 & $-0{,}9$\% \\
$\min A/B$ longitudinal & 2,273 & 2,36 & $-3{,}7$\% \\
$T_n$ [s] & 20,703 & 20,70 & 0,0\% \\
$\Hz$ [m] & 2,984 & 2,98 & +0,1\% \\
$\eta_{max}$ [m] & 11,346 & 11,35 & 0,0\% \\
\bottomrule
\end{tabular}
\end{table}

Os erros ficam entre 0,0\% e 3,7\%. O pequeno desvio no KG vem da altura do \emph{deck box}: com
$h_d = 7$~m, o KG seria de 32,99~m, praticamente o da aula, mas o $A/B$ subiria para 2,42 e 2,39 e
se afastaria dos valores de aula. Com $h_d = 8$~m os dois resultados ficam dentro de 4\%, e por
isso esse valor foi mantido. As curvas GZ e HA da plataforma de referência estão na
Figura~\ref{fig:valgz}.

\begin{figure}[!htb]
\centering
\includegraphics[width=\textwidth]{\RBA/f_gz_validacao.png}
\caption[Validação: curvas GZ e HA da plataforma de referência]{Validação: curvas GZ e HA da
plataforma dos Exercícios 7 a 10 nos dois planos.}\label{fig:valgz}
\end{figure}

\subsection{Verificações de consistência}

Além da comparação com as aulas, cinco verificações internas confirmam que as etapas do modelo são
coerentes entre si. A inclinação da curva GZ em $0^\circ$ é igual ao GM da hidrostática nos dois
planos (5,825~m, erro nulo), o que mostra que as projeções 2D reproduzem o Exercício 7. O RAO tende
a 1,000 quando $\omega \to 0$, como esperado para uma plataforma que acompanha a onda longa. O
$H_s$ recuperado da integral do espectro é de 12,21~m, com erro de 0,1\% devido à normalização
aproximada $1 - 0{,}287\ln\gamma$. A soma dos pesos fecha com o deslocamento, com erro nulo. Por
fim, a linha d'água inclinada obtida por Newton foi comparada à de bissecção em 2.030 geometrias,
nos dois planos e de 0 a $85^\circ$, com diferença máxima de $1{,}2\times10^{-8}$~m em GZ.

\subsection{Semissubmersíveis reais do pôster de 2011}\label{sec:val-poster}

A terceira frente usa o pôster \emph{2011 Worldwide Survey of Semi-FPSs and FPUs}~\cite{poster},
referência [6] do enunciado, que traz dimensões de \emph{pontoons} e colunas, calado e
deslocamento de unidades reais. Três delas têm a topologia do enunciado, com dois \emph{pontoons}
e quatro colunas. O modelo foi aplicado a elas usando o comprimento do \emph{pontoon} como $L_p$,
a largura e a altura do \emph{pontoon}, o lado da coluna e o calado máximo de operação, convertidos
de pés para metros (Tabela~\ref{tab:valposter}).

\begin{table}[!htb]
\centering
\caption[Validação do deslocamento com semissubmersíveis reais]{Deslocamento calculado pelo modelo
para semissubmersíveis reais de 2 \emph{pontoons} e 4 colunas. Dimensões do pôster de 2011, em
metros.}\label{tab:valposter}
\small\setlength{\tabcolsep}{4.5pt}
\begin{tabular}{lrrrrrrrrr}
\toprule
Unidade & $L_p$ & $b_p$ & $h_p$ & $b_c$ & $T$ & $\Delta$ pôster [t] & $\Delta$ modelo [t] & Erro & $T_n$ [s] \\
\midrule
Veslefrikk B (Noruega) & 107,6 & 14,3 & 9,5 & 12,5 & 25,9 & 43.305 & 40.395 & $-6{,}7$\% & 22,3 \\
Njord A (Noruega) & 100,0 & 17,7 & 9,1 & 15,9 & 23,2 & 45.077 & 47.571 & $+5{,}5$\% & 19,2 \\
P-20 (Brasil) & 80,5 & 15,9 & 7,6 & 12,8 & 20,4 & 25.983 & 28.524 & $+9{,}8$\% & 18,7 \\
\bottomrule
\end{tabular}
\end{table}

O modelo de caixas retangulares reproduz o deslocamento real com erro de até 10\%, com desvios
nos dois sentidos, que vêm de cantos arredondados, contraventamentos e do calado informado, que é o
máximo de operação. O período natural calculado fica entre 18,7 e 22,3~s, na mesma faixa do
critério $T_n > 20$~s. Essa validação se limita à hidrostática e ao período natural, porque o
pôster não traz dados de estabilidade nem de resposta.

\begin{sintese}
O modelo reproduz os resultados das aulas com erros de 0,0\% a 3,7\%, passa em todas as
verificações internas e estima o deslocamento de semissubmersíveis reais com erro de até 10\%. Ele
pode, portanto, ser usado com confiança na busca sistemática.
\end{sintese}

% =============================================================================
\section{Procedimento de variação sistemática e otimização}\label{sec:procedimento}

Com o modelo validado, o passo seguinte é explorar o domínio das variáveis e encontrar as
soluções que melhor equilibram peso e \emph{heave}. Esta seção apresenta a formulação do problema
e as quatro etapas do procedimento, sempre com a teoria do método seguida da forma como ele foi
aplicado.

\subsection{Formulação multiobjetivo}

O problema é
\begin{equation}
\min_{x}\ \big(\Wc(x),\ \Hz(x)\big) \quad \text{sujeito a} \quad g_j(x) \ge 0,\ j = 1,\ldots,m,
\qquad x_{lo} \le x \le x_{hi} ,
\end{equation}
com as restrições da Tabela~\ref{tab:restricoes} e o domínio da Tabela~\ref{tab:limites}. Diz-se
que uma solução $x^a$ \emph{domina} outra, $x^b$, quando não é pior em nenhum objetivo e é melhor
em pelo menos um, isto é, $\Wc^a \le \Wc^b$ e $\Hz^a \le \Hz^b$, com uma das desigualdades
estrita. As soluções viáveis que nenhuma outra domina formam a fronteira de Pareto; sobre ela,
reduzir um objetivo sempre piora o outro. O ponto utópico, com o menor peso e o menor $\Hz$ ao
mesmo tempo, não é viável.

Para obter essa fronteira, o procedimento combina as quatro etapas da Tabela~\ref{tab:etapas}:
uma amostragem ampla do domínio, que dá boas partidas; a soma ponderada e o
$\varepsilon$-restrito, que levam essas partidas até a fronteira; e um filtro de não dominância,
que separa os pontos ótimos.

\begin{table}[!htb]
\centering
\caption[Etapas da variação sistemática no caso-base]{Etapas da variação sistemática no
caso-base. Tempos com 8 processos em paralelo.}\label{tab:etapas}
\small
\begin{tabular}{>{\raggedright\arraybackslash}p{2.9cm}>{\raggedright\arraybackslash}p{6.1cm}>{\raggedright\arraybackslash}p{3.6cm}r}
\toprule
Etapa & O que faz & Resultado & Tempo \\
\midrule
1. Hipercubo latino & 20.000 combinações no domínio, avaliadas em dois estágios &
547 pré-viáveis, 227 viáveis (1,1\%) & 22 s \\
2. Soma ponderada & SLSQP para $w_1 = 0; 0{,}05; \ldots; 1$ (21 valores), até 6 partidas cada &
21 ótimos, 7 dominados & 182 s \\
3. $\varepsilon$-restrito & $\min \Wc$ com $\Hz \le \varepsilon$, 12 níveis, até 2 partidas cada &
12 ótimos & 35 s \\
4. Não dominância & filtra as 260 soluções viáveis e reconfirma cada candidato com GZ a cada
$0{,}25^\circ$ & 25 pontos (13 SP, 12 $\varepsilon$) & -- \\
\midrule
Rodada completa & validação, etapas 1 a 4, cenários e gráficos & & 241 s \\
\bottomrule
\end{tabular}
\end{table}

\subsection{Amostragem por hipercubo latino}

No hipercubo latino (LHS, McKay \emph{et al.}~\cite{mckay}), o intervalo de cada variável é
dividido em $n$ faixas de mesma probabilidade, e cada faixa recebe exatamente uma amostra; as faixas
das seis variáveis são combinadas por permutações aleatórias. Com o mesmo número de pontos, o LHS
cobre as distribuições marginais melhor que a amostragem aleatória simples, porque não deixa faixas
vazias nem concentra pontos.

Foram geradas 20.000 combinações, com semente 2026. Para economizar tempo, a avaliação é feita em
dois estágios. Primeiro são verificadas as restrições que não dependem da curva GZ (área, $GM_T$,
$T_n$, lastro, $b_c \le b_p$, \emph{air gap} e imersão das colunas), a 0,7~ms por amostra; só as
aprovadas passam pela curva GZ e pelo critério de áreas, a 6~ms. Das 20.000 amostras, 547
atendem ao primeiro estágio, mas só 227 (1,1\%) atendem também ao critério de áreas. O vento
reprova sozinho 320 cascos e é, de longe, a restrição mais seletiva. Os 227 cascos viáveis servem
de partida para as etapas seguintes.

\notacodigo{\texttt{variacao\_sistematica()} e \texttt{nao\_dominados()} (\lst{lhs}).}

\subsection{Otimizador SLSQP}

As duas etapas de otimização usam o SLSQP (\emph{Sequential Least Squares Programming},
Kraft~\cite{kraft}), um método de programação quadrática sequencial. A cada iteração, ele resolve um
subproblema quadrático que aproxima o lagrangiano, com hessiana atualizada pelo método BFGS e
restrições linearizadas, e ajusta o passo por busca linear em uma função de mérito. O método
converge para um ponto que satisfaz as condições de Karush-Kuhn-Tucker, ou seja, para um ótimo
local.

Foi usada a implementação do SciPy~\cite{scipy}, com até 200 iterações, $\texttt{ftol} = 10^{-8}$
e gradientes por diferenças finitas com passo de $10^{-5}$. Por ser um método local, cada problema
parte de vários pontos (\emph{multi-start}), e o resultado de um peso serve de partida para o
seguinte (continuação). Objetivo e restrições compartilham as avaliações por meio de um
\emph{cache}. Durante a busca, a GZ usa passo de $1^\circ$; ao final, todo ótimo é reavaliado com
passo de $0{,}25^\circ$ e só é aceito se continuar viável. Um resultado com
\texttt{success = False}, que indica limite de iterações ou busca linear sem progresso, é aceito
se for viável, porque o que importa para a fronteira é a viabilidade e o valor dos objetivos.

\notacodigo{\emph{cache} do problema e partida do SLSQP em \texttt{Problema} e
\texttt{\_tarefa\_slsqp()} (\lst{problema}).}

\subsection{Soma ponderada}

A soma ponderada transforma os dois objetivos em um só,
\begin{equation}
F = w_1\, \Hz + w_2\, \frac{\Wc}{10.000~\text{t}}, \qquad w_1 + w_2 = 1 ,
\end{equation}
em que a divisão por 10.000~t coloca os dois termos na mesma ordem de grandeza (1 a 2,4). Ao variar
$w_1$ de 0 a 1, o mínimo de $F$ percorre a fronteira, mas apenas nos trechos convexos: esse mínimo é
sempre um ponto do fecho convexo da fronteira, e um trecho côncavo não é atingido por nenhum peso.
Os extremos $w_1 = 0$ e $w_1 = 1$ dão as soluções de mínimo peso e de mínimo \emph{heave}.

Para cada um dos 21 valores de $w_1$ (passo 0,05), o SLSQP partiu das 5 amostras do LHS com menor
$F$ e do ótimo do $w_1$ anterior, com as partidas rodando em paralelo. O mínimo peso obtido, de
10.346~t com $w_1 = 0$, é 14,9\% mais leve que a melhor amostra viável do LHS. Por outro lado,
sete dos 21 ótimos ($w_1$ = 0,05; 0,15; 0,25; 0,40; 0,50; 0,55 e 0,60) ficaram dominados por
outros. Isso mostra que a fronteira não é convexa: nesses trechos, o SLSQP para em ótimos locais,
vários deles com \texttt{success = False}.

\notacodigo{\texttt{soma\_ponderada()} (\lst{somapond}).}

\subsection{Método $\varepsilon$-restrito}

Para alcançar os trechos que a soma ponderada não atinge, o método $\varepsilon$-restrito (Haimes
\emph{et al.}~\cite{haimes}) minimiza um objetivo e transforma o outro em restrição:
\begin{equation}
\min_x \Wc(x) \quad \text{sujeito a} \quad \Hz(x) \le \varepsilon \quad \text{e} \quad g_j(x) \ge 0 .
\end{equation}
Variando $\varepsilon$, qualquer ponto da fronteira pode ser obtido, inclusive nos trechos não
convexos~\cite{miettinen}.

Foram usados 12 níveis de $\varepsilon$ igualmente espaçados entre o menor e o maior $\Hz$ da soma
ponderada (1,0435 e 2,3664~m), excluídos os extremos, com passo de 0,102~m. Cada nível partiu da
solução conhecida mais leve com $\Hz \le \varepsilon$ e do ótimo da soma ponderada logo abaixo de
$\varepsilon$. Os 12 níveis terminaram em pontos viáveis, todos com a restrição
$\Hz \le \varepsilon$ ativa.

\notacodigo{\texttt{epsilon\_restrito()} e \texttt{ponto\_joelho()} (\lst{epsilon}).}

\subsection{Filtro de não dominância}

As 260 soluções viáveis (227 do LHS, 21 da soma ponderada e 12 do $\varepsilon$-restrito) são
ordenadas por peso crescente, e cada uma é mantida se o seu $\Hz$ for menor que o menor $\Hz$ das
anteriores, com tolerância relativa de $10^{-5}$; o algoritmo é $O(n\log n)$. Os candidatos do LHS
que entram na fronteira são reavaliados com GZ a cada $0{,}25^\circ$, e os que falham são
descartados antes de repetir o filtro. A fronteira final tem 25 pontos, 13 da soma ponderada e 12
do $\varepsilon$-restrito. Nenhuma amostra do LHS sobrevive, o que mostra o ganho da etapa de
otimização sobre a amostragem.

\notacodigo{montagem da fronteira no bloco principal (\lst{fronteira}).}

\subsection{Escolha do ponto de compromisso}\label{sec:joelho}

Resta escolher uma solução sobre a fronteira. Para isso, os dois objetivos são normalizados entre
0 e 1 ao longo da fronteira,
\begin{equation}
\hat W = \frac{\Wc - \Wc^{\text{mín}}}{\Wc^{\text{máx}} - \Wc^{\text{mín}}}, \qquad
\hat H = \frac{\Hz - \Hz^{\text{mín}}}{\Hz^{\text{máx}} - \Hz^{\text{mín}}},
\end{equation}
e a solução intermediária é o ponto da fronteira mais próximo do ponto utópico $(0, 0)$, o
\emph{joelho} da curva. A Seção~\ref{sec:escolha} mostra que essa escolha coincide com a mudança
de patamar do custo marginal de reduzir o \emph{heave}.


% =============================================================================
\section{Conjunto viável e fronteira de Pareto}\label{sec:pareto}

O procedimento da seção anterior produz dois resultados: o mapa do conjunto viável, que mostra
onde estão as plataformas possíveis, e a fronteira de Pareto, que contém as melhores entre elas.
Esta seção apresenta os dois e explica por que a fronteira tem a forma que tem.

\subsection{Conjunto viável}

A Figura~\ref{fig:pareto} mostra as 20.000 amostras no plano peso $\times$ $\Hz$, separadas em
inviáveis, reprovadas só no critério de vento e viáveis, junto com os ótimos da soma ponderada e do
$\varepsilon$-restrito e a fronteira de Pareto.

\begin{figure}[!htb]
\centering
\includegraphics[width=0.9\textwidth]{\RBA/f_pareto.png}
\caption[Conjunto viável e fronteira de Pareto (caso-base)]{Conjunto viável e fronteira de Pareto
no plano peso estrutural $\times$ $\Hz$ (caso-base).}\label{fig:pareto}
\end{figure}

A figura deixa claro que o vento é o gargalo do projeto. O critério de áreas reprova sozinho 320 dos
547 cascos que atendem às demais restrições, e essas amostras, em laranja, ocupam justamente a
região de baixo peso onde a fronteira estaria se não houvesse o critério. A figura também mostra o
ganho do otimizador: nenhuma amostra do LHS sobrevive ao filtro de não dominância, e o SLSQP
reduz o peso mínimo em 14,9\% em relação à melhor amostra viável. A amostragem serve, portanto,
para dar boas partidas, mas não substitui a otimização.

\subsection{Fronteira de Pareto}

A fronteira tem 25 pontos, de 10.346~t ($\Hz = 2{,}37$~m) a 23.634~t ($\Hz = 1{,}04$~m), listados no
Apêndice~\ref{ap:fronteira}. Ela é não convexa, já que 7 dos 21 pontos da soma ponderada ficam
dominados e o $\varepsilon$-restrito precisa preencher os vazios. Entre os dois extremos, cortar o
\emph{heave} em 56\% custa 128\% a mais de peso. A solução escolhida, discutida na
Seção~\ref{sec:escolha}, captura 62\% do ganho de movimento com apenas 35\% do acréscimo de massa.

\subsection{Restrições ativas: um projeto governado pela estabilidade}

A Tabela~\ref{tab:ativasbase} conta, ao longo dos 25 pontos da fronteira, quantas vezes cada
restrição fica ativa. O $GM_T$ no máximo de 4~m e o $A/B$ no mínimo de 1,30 aparecem em todos os
pontos.

\begin{table}[!htb]
\centering
\caption{Restrições ativas ao longo dos 25 pontos da fronteira (caso-base).}\label{tab:ativasbase}
\begin{tabular}{lr}
\toprule
Restrição ativa & Pontos \\
\midrule
$GM_T = 4$~m & 25 / 25 \\
$A/B = 1{,}30$ & 25 / 25 \\
lastro positivo (11 t) & 17 / 25 \\
$b_c = b_p$ & 9 / 25 \\
área de convés mínima & 4 / 25 \\
\emph{Air gap} e $T_n$ & 0 / 25 \\
\bottomrule
\end{tabular}
\end{table}

O mecanismo por trás desse resultado é o seguinte. O \emph{topside} de 22.420~t fica a cerca de
60~m da quilha e eleva o KG. Para atender ao critério de áreas com o menor casco possível, o
otimizador leva o GM ao máximo permitido, porque isso aumenta a GZ em todos os ângulos. Com o GM
travado em 4~m, a área $A$ só pode crescer adiando o alagamento, o que pede colunas mais altas e,
com elas, mais peso e mais área exposta ao vento. Já o \emph{air gap} e o $T_n$ nunca ficam ativos:
o mar do campo é brando, e a altura das colunas é definida pelo ângulo de alagamento, e não pela
crista da onda.

\subsection{Duas famílias de projeto}

A Figura~\ref{fig:variaveis} mostra as variáveis ao longo da fronteira, normalizadas pelos limites
do domínio. Um salto entre 13.657~t e 14.241~t separa duas famílias de projeto. Até cerca de
13.700~t, a fronteira é formada por cascos rasos e largos, com $T \approx 12$--19~m, $b_c$ pequeno e
$L_p$ crescendo até 150~m; a área de linha d'água pequena mantém $T_n$ entre 29 e 40~s, e o
\emph{heave} cai com o aumento de $h_c$ e $L_p$. Acima de cerca de 14.200~t, os cascos passam a ser
compactos e profundos, com $L_p$ perto do mínimo, colunas maiores e $T$ de 23 a 40~m; aqui o
\emph{heave} cai porque os \emph{pontoons} ficam mais fundos e a excitação se atenua, com $T_n$
entre 21,6 e 23~s.

\begin{figure}[!htb]
\centering
\includegraphics[width=0.85\textwidth]{\RBA/f_variaveis_fronteira.png}
\caption[Variáveis de projeto ao longo da fronteira (caso-base)]{Variáveis de projeto normalizadas
pelos limites ao longo da fronteira (0 = mínimo, 1 = máximo).}\label{fig:variaveis}
\end{figure}

\subsection{Vista lateral das três soluções}

A Figura~\ref{fig:geo3} desenha em escala os dois extremos da fronteira e a solução escolhida, com
a linha d'água de projeto e a crista máxima. A de mínimo peso é rasa e larga; a de mínimo
\emph{heave} tem calado e colunas máximos; e a escolhida fica entre as duas, já na família
compacta.

\begin{figure}[!htb]
\centering
\includegraphics[width=\textwidth]{\RB/05_geometria_3_cenarios.png}
\caption[Vista lateral dos três cenários (caso-base)]{Vista lateral dos três cenários (caso-base),
com a linha d'água de projeto e a crista máxima; convés e módulos em cinza.}\label{fig:geo3}
\end{figure}

% =============================================================================
\section{Escolha da solução intermediária}\label{sec:escolha}

A solução intermediária é o joelho da fronteira (Seção~\ref{sec:joelho}), obtido pelo
$\varepsilon$-restrito com $\varepsilon = 1{,}552$~m. Mais do que uma regra geométrica, essa escolha
se apoia em quatro argumentos de projeto.

O primeiro é o custo marginal de reduzir o \emph{heave}, que muda de patamar no joelho
(Tabela~\ref{tab:custo} e Figura~\ref{fig:custo}). Até a solução escolhida, cada milímetro de
$\Hz$ custa 5,7~t de aço; depois dela, custa 17,0~t, três vezes mais.

\begin{table}[!htb]
\centering
\caption{Custo marginal de redução do \emph{heave} ao longo da fronteira (caso-base).}\label{tab:custo}
\begin{tabular}{lrrr}
\toprule
Trecho & Peso adicional [t] & Redução de $\Hz$ [m] & Custo [t/mm] \\
\midrule
Cenário 1 $\to$ 3 & 4.652 & 0,81 & 5,7 \\
Cenário 3 $\to$ 2 & 8.635 & 0,51 & 17,0 \\
\bottomrule
\end{tabular}
\end{table}

\begin{figure}[!htb]
\centering
\includegraphics[width=0.72\textwidth]{\RBA/f_custo_marginal.png}
\caption[Custo marginal de reduzir o \emph{heave}]{Custo marginal de reduzir o \emph{heave}:
5,7~t por mm até a solução escolhida e 17,0~t por mm depois dela.}\label{fig:custo}
\end{figure}

O segundo argumento é o equilíbrio entre os extremos. Em relação ao cenário de mínimo peso, a
solução escolhida corta 34,4\% do $\Hz$ com 45,0\% a mais de casco; em relação ao de mínimo
\emph{heave}, economiza 36,5\% do aço, com um $\Hz$ 48,8\% maior. O terceiro é a folga geométrica:
das três soluções, é a única com todas as variáveis no interior do domínio, enquanto os extremos
encostam nos limites ($b_c$, $h_c$ e $T$ mínimos no cenário 1; $h_c$ e $T$ máximos no cenário 2) e
dependem, portanto, de escolhas de domínio que não vêm do enunciado. O quarto é a reserva de
estabilidade: o ângulo de alagamento é de 21,6$^\circ$, mais que o dobro dos 9,0$^\circ$ do casco
de mínimo peso.

O ponto fraco da solução escolhida é a ausência de lastro, com apenas 11~t de margem numérica. Essa
questão é tratada na Seção~\ref{sec:lastro}, que repete o projeto com lastro mínimo obrigatório.

% =============================================================================
\section{Ficha técnica comparativa dos três cenários}\label{sec:cenarios}

O enunciado pede a comparação de três cenários: os dois extremos da fronteira ($w_1 = 0$ e
$w_1 = 1$) e a solução escolhida no joelho. A Tabela~\ref{tab:ficha} reúne suas características,
e a Figura~\ref{fig:heave3} compara a resposta em \emph{heave} dos três.

\needspace{0.75\textheight}
{\small\renewcommand{\arraystretch}{1.04}
\begin{longtable}{lrrr}
\caption{Ficha técnica comparativa dos três cenários (caso-base).}\label{tab:ficha}\\
\toprule
 & 1. Mín. peso & 2. Mín. \emph{heave} & 3. Escolhida \\
\midrule
\endfirsthead
\toprule
 & 1. Mín. peso & 2. Mín. \emph{heave} & 3. Escolhida \\
\midrule
\endhead
Origem & $w_1 = 0$ & $w_1 = 1$ & $\varepsilon = 1{,}552$ m \\
$L_p$ [m] & 137,49 & 110,11 & 111,80 \\
$b_p$ [m] & 16,39 & 14,22 & 13,40 \\
$h_p$ [m] & 7,37 & 15,03 & 10,21 \\
$b_c$ [m] & 8,00 & 14,22 & 12,55 \\
$h_c$ [m] & 20,00 & 50,00 & 39,61 \\
$T$ [m] & 12,00 & 40,00 & 23,98 \\
Área de convés [m$^2$] & 18.905 & 12.125 & 12.499 \\
\midrule
$\nabla$ [m$^3$] & 34.384 & 67.280 & 39.265 \\
$\Delta$ [t] & 35.244 & 68.962 & 40.247 \\
KB [m] & 3,89 & 13,52 & 7,75 \\
$BM_T$ / $BM_L$ [m] & 27,34 / 31,25 & 27,85 / 27,85 & 39,06 / 39,74 \\
$KM_T$ / $KM_L$ [m] & 31,23 / 35,14 & 41,37 / 41,37 & 46,82 / 47,49 \\
\midrule
\textbf{Peso do casco [t]} & \textbf{10.346} & \textbf{23.634} & \textbf{14.998} \\
\quad \emph{pontoons} / colunas [t] & 8.963 / 1.382 & 12.710 / 10.924 & 8.257 / 6.741 \\
Peso da planta [t] & 13.188 & 13.188 & 13.188 \\
Peso do \emph{topside} [t] & 22.420 & 22.420 & 22.420 \\
\emph{Risers} + amarração [t] & 2.467 & 4.827 & 2.817 \\
Lastro [t] & 11 & 18.081 & 11 \\
$V_{lastro}/V_{pontoons}$ & 0,000 & 0,375 & 0,000 \\
KG [m] & 27,23 & 37,37 & 42,82 \\
$GM_T$ / $GM_L$ [m] & 4,00 / 7,91 & 4,00 / 4,00 & 4,00 / 4,67 \\
\midrule
$\phial$ transversal / longitudinal [$^\circ$] & 8,96 / 9,77 & 24,45 / 24,45 & 21,57 / 23,86 \\
$\min A/B$ (transversal / longitudinal) & 1,302 / 1,601 & 1,304 / 1,486 & 1,303 / 1,304 \\
\midrule
$A_{33}$ [t] & 70.540 & 40.568 & 36.378 \\
$C_{33}$ [kN/m] & 2.574 & 8.136 & 6.337 \\
$T_n$ [s] & 40,28 & 23,05 & 21,85 \\
\textbf{$\Hz$ com $\gamma = 3{,}3$ [m]} & \textbf{2,366} & \textbf{1,043} & \textbf{1,552} \\
$\Hz$ com $\gamma = 1$ [m] & 2,413 & 0,993 & 1,454 \\
$\Hz / H_s$ & 42\% & 19\% & 28\% \\
\midrule
$\eta_{max}$ / $z_{min}$ [m] & 5,21 / $-2{,}20$ & 5,21 / $-0{,}97$ & 5,21 / $-1{,}44$ \\
$\amin$ / $a$ [m] & 7,41 / 15,37 & 6,18 / 25,03 & 6,65 / 25,84 \\
\midrule
Restrições ativas & $GM_T$, lastro, $A/B$ & $GM_T$, $b_c{=}b_p$, $A/B$ & $GM_T$, lastro, $A/B$ \\
Variáveis no limite & $b_c$, $h_c$, $T$ mín. & $h_c$, $T$ máx. & nenhuma \\
\bottomrule
\end{longtable}}

\begin{figure}[!htb]
\centering
\includegraphics[width=0.7\textwidth]{\RBA/f_heave_3cenarios.png}
\caption[RAO e espectro de resposta dos três cenários (caso-base)]{RAO e espectro de resposta em
\emph{heave} dos três cenários (caso-base); em azul claro, o espectro do mar.}\label{fig:heave3}
\end{figure}

\subsection{Cenário 1: mínimo peso estrutural}

O cenário de mínimo peso é uma semissubmersível rasa: $T$, $h_c$ e $b_c$ estão no limite mínimo, e
os \emph{pontoons} levam 87\% do peso do casco. Com colunas de $8\times8$~m, o GM só é obtido pelo
afastamento entre elas, o que leva a $L_p = 137{,}5$~m e a 63\% de folga na área de convés. A área
de linha d'água de 256~m$^2$ dá $T_n = 40$~s, longe da ressonância, e ainda assim o \emph{heave} é
o maior dos três ($\Hz/H_s = 42\%$). A razão está na profundidade dos \emph{pontoons}: o topo fica a
só 4,6~m da superfície, e o fator de atenuação $e^{-k(T-h_p)}$ vale 0,92 no pico do mar, de modo que
a força de Froude-Krylov quase não se atenua (Figura~\ref{fig:heave3}).

A estabilidade também é frágil. O ângulo de alagamento é de 9,0$^\circ$, a GZ fica negativa a partir
de 17,8$^\circ$ (Figura~\ref{fig:gz1}), e o vento de tempestade já inclina a plataforma 4,9$^\circ$
em equilíbrio, mais da metade do caminho até o alagamento. O critério é atendido, mas a reserva de
estabilidade é pequena demais para um projeto real.

\begin{figure}[!htb]
\centering
\includegraphics[width=\textwidth]{\RBA/f_gz_cen1.png}
\caption{Cenário 1 (mínimo peso): curvas GZ e HA.}\label{fig:gz1}
\end{figure}

\subsection{Cenário 2: mínimo \emph{heave}}

No extremo oposto, o cenário de mínimo \emph{heave} é uma semissubmersível de calado profundo, com
$T = 40$~m e $h_c = 50$~m no máximo e $b_c = b_p = 14{,}2$~m. Com o topo dos \emph{pontoons} a 25~m
de profundidade, o fator de atenuação cai para 0,65 no pico do mar, e $\Hz/H_s$ fica em 19\%. É o
único cenário com lastro de fato, 18.081~t ou 37,5\% do volume dos \emph{pontoons}, o que dá margem
operacional. Por ter $b_c = b_p$ e planta quadrada, tem $GM_L = GM_T$ e $\phial = 24{,}4^\circ$
nos dois planos; ainda assim, o $A/B$ é de 1,30 no plano transversal e de 1,49 no longitudinal,
por causa da diferença de $z_{CLR}$ (13,5 e 20,0~m). O preço é um casco 2,3 vezes mais pesado que o
do cenário 1, com 46\% do peso nas colunas, e o final da fronteira é quase plano: de 23.327~t para
23.634~t, o $\Hz$ cai só de 1,046 para 1,043~m (Figura~\ref{fig:gz2}).

\begin{figure}[!htb]
\centering
\includegraphics[width=\textwidth]{\RBA/f_gz_cen2.png}
\caption{Cenário 2 (mínimo \emph{heave}): curvas GZ e HA.}\label{fig:gz2}
\end{figure}

\subsection{Cenário 3: solução escolhida}

A solução escolhida combina características dos dois extremos: calado intermediário (24~m),
colunas de 12,6~m e \emph{pontoons} com topo a 13,8~m de profundidade, onde o fator de atenuação
vale 0,79 no pico do mar. O resultado é $\Hz/H_s = 28\%$ com $T_n = 21{,}85$~s. A
Seção~\ref{sec:final} detalha essa solução.

% =============================================================================
\section{Comparativo com plataformas reais}\label{sec:reais}

O modelo paramétrico usa correlações simples de peso e área, e por isso vale perguntar se as
soluções que ele produz se parecem com plataformas que existem. Esta seção compara a solução
escolhida com a frota de semissubmersíveis do pôster de 2011, com as unidades que operam no Baixo
Congo e com semissubmersíveis de produção recentes de porte semelhante.

\subsection{Semissubmersíveis do pôster de 2011}

O pôster de 2011~\cite{poster} lista 50 unidades. Os valores usados foram lidos da tabela do pôster e
convertidos de pés para metros; como a extração de texto desalinha algumas linhas, só entram as
unidades cuja sequência de campos é inequívoca, registradas em
\path{caso_base/resultados/complementar/poster2011_unidades.csv}. A validação do casco com três
dessas unidades foi apresentada na Seção~\ref{sec:val-poster}. Aqui, a Tabela~\ref{tab:frota} e a
Figura~\ref{fig:poster} comparam a solução escolhida com a frota.

\begin{table}[!htb]
\centering
\caption[Solução escolhida contra a frota do pôster de 2011]{Solução escolhida contra as
semissubmersíveis do pôster de 2011 ($n$ = número de unidades com o dado).}\label{tab:frota}
\small
\begin{tabular}{lrrrrrr}
\toprule
Grandeza & $n$ & Mínimo & Mediana & Máximo & G02 & G02 (10\%) \\
\midrule
Calado $T$ [m] & 10 & 20,4 & 25,9 & 32,0 & 24,0 & 25,9 \\
Altura da coluna $h_c$ [m] & 7 & 33,2 & 39,9 & 43,9 & 39,6 & 36,5 \\
Lado da coluna $b_c$ [m] & 11 & 12,2 & 15,8 & 19,8 & 12,6 & 12,9 \\
Largura do \emph{pontoon} $b_p$ [m] & 11 & 9,1 & 15,8 & 19,8 & 13,4 & 13,2 \\
Altura do \emph{pontoon} $h_p$ [m] & 11 & 6,1 & 8,5 & 10,1 & 10,2 & 12,4 \\
\emph{Topside} por mil bopd [t] & 6 & 100 & 145 & 252 & 141 & 141 \\
$\Delta$ por mil bopd [t] & 7 & 286 & 526 & 889 & 253 & 285 \\
\bottomrule
\end{tabular}
\end{table}

\begin{figure}[!htb]
\centering
\includegraphics[width=\textwidth]{\RBA/f_poster2011.png}
\caption[Solução escolhida contra o pôster de 2011]{Solução escolhida contra as
semissubmersíveis do pôster de 2011: (a) dimensões do casco; (b) peso de \emph{topside} por mil
bopd; (c) deslocamento por mil bopd.}\label{fig:poster}
\end{figure}

Nas dimensões, a solução se encaixa bem na frota. Calado, altura e lado da coluna e largura do
\emph{pontoon} ficam dentro da faixa, e a altura da coluna (39,6~m) praticamente coincide com a
mediana (39,9~m). As colunas, de 12,6~m, estão no limite inferior da frota, o que é coerente com o
$GM_T$ ativo em 4~m, já que colunas maiores aumentariam o BM. A única dimensão fora da faixa é a
altura do \emph{pontoon}, um pouco acima do máximo do pôster (10,2 contra 10,1~m).

O \emph{topside} também confere. A correlação da disciplina, com planta de 83~kg/boe mais 35\% de
convés e 35\% de carga operacional, dá 141~t por mil bopd, praticamente a mediana do pôster
(145~t) e próxima de P-52 (134~t) e Atlantis (128~t). O deslocamento é onde aparece a diferença: a
solução escolhida desloca 253~t por mil bopd, abaixo da unidade mais enxuta do pôster, a Troll C,
com 286~t. Com lastro mínimo de 10\%, o valor sobe para 285~t, o nível da Troll C. O casco do
caso-base é leve, portanto, porque não reserva deslocamento para lastro, e não por erro de
hidrostática ou de \emph{topside}.

\subsection{O que opera no Baixo Congo}

As unidades de produção em águas profundas da bacia do Baixo Congo são, em sua maioria, FPSOs com
ancoragem distribuída (Tabela~\ref{tab:reaisregiao}). O análogo mais próximo do campo é o FPSO do
projeto CLOV~\cite{clov}, da TotalEnergies, no Bloco 17 de Angola, quase um gêmeo do KZ-2026-B em
capacidade (160 contra 159 mil bpd) e em lâmina d'água (1.291 contra 1.450~m), com casco de navio e
armazenamento de 1,8 milhão de barris. A Moho/Bilondo, da Total, em 620~m de lâmina d'água na
República do Congo, é a única unidade da região no pôster de 2011; aparece como FPU, mas também tem
casco de navio (618 × 112~pés). A única semissubmersível de produção da região é a
Scarabeo~5~\cite{scarabeo5}, uma sonda de 1990 convertida pela Saipem em unidade de tratamento de
gás para o projeto Congo LNG da Eni, em 35~m de lâmina d'água. Seu comprimento de 111~m coincide
com o $L_p = 111{,}8$~m da solução escolhida, mas os dados do casco não são públicos, e ela não
entrou na comparação numérica.

\begin{table}[!htb]
\centering
\caption[Campo KZ-2026-B e plataformas reais da região]{Campo KZ-2026-B e plataformas reais da
região. Dados públicos.}\label{tab:reaisregiao}
\small
\begin{tabular}{>{\raggedright\arraybackslash}p{2.7cm}>{\raggedright\arraybackslash}p{3.8cm}>{\raggedright\arraybackslash}p{3.8cm}>{\raggedright\arraybackslash}p{3.8cm}}
\toprule
 & KZ-2026-B (este projeto) & CLOV FPSO & Scarabeo 5 FPU \\
\midrule
Local & Baixo Congo, Angola & Bloco 17, Angola (Baixo Congo) & Congo LNG, Rep. do Congo \\
Distância à costa & 190 km & 140 km de Luanda & próxima ao terminal de Djeno \\
Lâmina d'água & 1.450 m & 1.291 m & 35 m \\
Produção & 158.894 bpd (pico) & 160.000 bpd de óleo & gás (separação e compressão) \\
Casco & semi, 2 \emph{pontoons} + 4 colunas & navio, 305 × 61 m, calado 24,4 m & semi, 2 \emph{pontoons} + 6 colunas, 111 × 73 m \\
Armazenamento & nenhum & 1,8 milhão de bbl & nenhum \\
Peso / deslocamento & $\Delta$ = 40.247 t & 110.000 t (peso informado) & -- \\
\bottomrule
\end{tabular}
\end{table}

Na interpretação do grupo, dois dados do campo explicam a preferência pelo FPSO na região. O
primeiro é o escoamento: o CLOV armazena cerca de 11 dias de produção e exporta o óleo por uma boia
de descarga a 1,8~km da unidade, enquanto a semissubmersível não tem armazenamento e exigiria um
oleoduto de 190~km até a costa ou um FSO ao lado, custo que o modelo paramétrico não inclui. O
segundo é o mar brando: com $H_s = 5{,}6$~m no mar de 100 anos, o movimento de um navio é
aceitável, e a vantagem da semissubmersível em \emph{heave} pesa mais em mares severos, como o do
Golfo do México, onde ela é a solução preferida~\cite{omsemis}.

\subsection{Semissubmersíveis de produção de porte semelhante}

Para comparar o casco com unidades recentes de capacidade e lâmina d'água semelhantes, foram usadas
semissubmersíveis de produção do Golfo do México~\cite{omsemis,omgom}. A mais próxima é a Argos,
da bp, com 140 mil boe/d em 1.372~m de lâmina d'água (Tabela~\ref{tab:reaissemi} e
Figura~\ref{fig:reais}).

\begin{table}[!htb]
\centering
\caption[Semissubmersíveis de produção reais e soluções do grupo]{Semissubmersíveis de produção
reais e soluções do grupo. Os pesos da Argos foram convertidos de toneladas curtas; o peso da King's
Quay é o peso da unidade (casco + \emph{topside}).}\label{tab:reaissemi}
\small
\begin{tabular}{lrrrrr}
\toprule
Unidade & Capacidade & $\Delta$ [t] & $\Delta$/capacidade & Peso [t] & Peso/$\Delta$ \\
 & [mil bpd] & & [t por mil bpd] & & \\
\midrule
Argos (Mad Dog 2) & 140 & 101.000 & 721 & $\approx$ 54.400 & 0,54 \\
King's Quay & 80 & 38.000 & 475 & 21.498 & 0,57 \\
Vito & 100 & 38.000 & 380 & -- & -- \\
Anchor & 75 & 69.000 & 920 & -- & -- \\
\midrule
G02 · cenário 2 (mín. \emph{heave}) & 159 & 68.962 & 434 & 46.054 & 0,67 \\
G02 · cenário 3 (escolhida) & 159 & 40.247 & 253 & 37.418 & 0,93 \\
G02 · escolhida com lastro $\ge$ 10\% & 159 & 45.223 & 285 & 38.446 & 0,85 \\
\bottomrule
\end{tabular}
\end{table}

\begin{figure}[!htb]
\centering
\includegraphics[width=\textwidth]{\RBA/f_comparativo_reais.png}
\caption[Comparativo com semissubmersíveis de produção reais]{Comparativo com semissubmersíveis de
produção reais: (a) deslocamento por mil bpd de capacidade; (b) área de convés contra a correlação
da disciplina; (c) fração do deslocamento ocupada por casco e \emph{topside}.}\label{fig:reais}
\end{figure}

A solução escolhida desloca 253~t por mil bpd, de 1,5 a 3,6 vezes menos que as semissubmersíveis
reais (380 a 920~t por mil bpd), enquanto o cenário 2, com 434~t por mil bpd, fica na faixa das
unidades enxutas recentes, Vito e King's Quay. Três fatores explicam a diferença. O principal é o
lastro: casco e \emph{topside} ocupam 93\% do deslocamento da solução escolhida, contra 54\% a 57\%
nas unidades reais, que reservam deslocamento para lastro, cargas variáveis, \emph{risers} e
amarração; com lastro mínimo de 10\%, a fração cai para 85\%. O segundo é o ambiente, já que as
unidades do Golfo são projetadas para furacões (a King's Quay, para a tempestade de 1.000
anos~\cite{omgom}), com mar muito mais severo que o de Angola, onde o \emph{air gap} tem folga de
19~m e o calado não é exigido pelo movimento. O terceiro são as próprias correlações de peso, que
somam \emph{risers} e amarração em 7\% de $\Delta$ e aplicam $\gamma_e = 0{,}270$~t/m$^3$ ao
volume envolvido, sem margens de projeto.

A área de convés mostra o efeito oposto. A correlação da disciplina leva a 11.563~m$^2$ para 159 mil
bpd, o equivalente a um único nível, enquanto a Argos tem convés em planta de cerca de 110 × 49~m
($\approx$ 5.400~m$^2$) e a Shenandoah, 91 × 91~m (8.281~m$^2$) para 100 mil bpd, ambas com vários
níveis. Por isso o $L_p$ do projeto é maior que o lado das unidades reais. Como a face inferior do
convés ($L_p^2$) gera a maior parcela do momento de vento, a correlação de um nível é conservadora
para o critério de áreas. Ainda assim, a altura total da solução escolhida, de 69,8~m da quilha ao
topo dos módulos, fica abaixo dos 90~m da Shenandoah.

\begin{sintese}
O modelo reproduz a ordem de grandeza e a configuração de semissubmersíveis reais e estima bem o
\emph{topside}, mas produz cascos mais leves porque não reserva lastro e porque o mar do campo é
brando; com lastro mínimo de 10\%, o deslocamento específico chega ao da unidade mais enxuta do
pôster. No Baixo Congo, os análogos reais têm casco de navio, e uma semissubmersível só seria
competitiva se houvesse infraestrutura de exportação; o custo de um FSO ou de um oleoduto de
190~km deve entrar na comparação de conceitos.
\end{sintese}


% =============================================================================
\section{Solução final}\label{sec:final}

Esta seção detalha a solução escolhida (cenário 3) na mesma sequência do modelo: dimensões e pesos,
estabilidade e critério de vento, \emph{heave}, \emph{air gap} e, por fim, o quadro completo das
restrições.

\subsection{Dimensões, pesos e centros}

A solução escolhida tem $L_p = 111{,}80$~m, $b_p = 13{,}40$~m, $h_p = 10{,}21$~m, $b_c = 12{,}55$~m,
$h_c = 39{,}61$~m e $T = 23{,}98$~m, com \emph{air gap} de 25,84~m. A área de convés é de
12.499~m$^2$, 8,1\% acima da mínima, e todas as variáveis estão no interior do domínio. A
Tabela~\ref{tab:pesosfinal} mostra a composição do peso.

\begin{table}[!htb]
\centering
\caption{Pesos e centros da solução final (caso-base).}\label{tab:pesosfinal}
\begin{tabular}{lrrr}
\toprule
Item & Peso [t] & VCG [m] & Momento [t$\cdot$m] \\
\midrule
Planta de produção & 13.188,2 & 63,82 & 841.656 \\
Estrutura do convés & 4.615,9 & 53,82 & 248.421 \\
Carga operacional & 4.615,9 & 53,82 & 248.421 \\
\emph{Pontoons} & 8.257,3 & 5,10 & 42.136 \\
Colunas & 6.741,0 & 30,01 & 202.312 \\
\emph{Risers} & 1.609,9 & 49,82 & 80.201 \\
Amarração & 1.207,4 & 49,82 & 60.151 \\
Lastro & 11,0 & 0,00 & 0 \\
\midrule
\textbf{Total ($=\Delta$)} & \textbf{40.246,6} & \textbf{42,82} & \textbf{1.723.298} \\
\bottomrule
\end{tabular}
\end{table}

O casco pesa 14.998~t, e o \emph{topside}, 22.420~t, com centro a 59,7~m da quilha. É esse
\emph{topside} alto que eleva o KG a 42,82~m. Com KB = 7,75~m e $KM_T = 46{,}82$~m, resultam
$GM_T = 4{,}00$~m e $GM_L = 4{,}67$~m: o GM de 4~m vem do BM de 39~m, isto é, do afastamento entre
as colunas.

\subsection{Estabilidade intacta e critério de vento}

A Figura~\ref{fig:gzfinal} mostra as curvas GZ e HA nos dois planos. As forças e os braços de cada
parcela de vento estão na Tabela~\ref{tab:ventofinal}, os momentos na Tabela~\ref{tab:momentos} e
o resultado do critério de áreas nos quatro casos na Tabela~\ref{tab:criterio}.

\begin{figure}[!htb]
\centering
\includegraphics[width=\textwidth]{\RBA/f_gz_final.png}
\caption[Solução final: curvas GZ e HA]{Solução final: curvas GZ e HA nos planos transversal e
longitudinal, com $\phial$ e os interceptos ($\bullet$ $\varphi_1$, $\blacksquare$
$\varphi_2$).}\label{fig:gzfinal}
\end{figure}

\begin{table}[!htb]
\centering
\caption[Forças e braços de vento na solução final]{Forças e braços de vento por parcela na solução
final. As forças são iguais nos dois planos (planta quadrada); os braços dependem do $z_{CLR}$ de
cada plano.}\label{tab:ventofinal}
\small\setlength{\tabcolsep}{4.5pt}
\begin{tabular}{lrrrrrrrr}
\toprule
 & & & & & \multicolumn{2}{c}{Força [kN]} & \multicolumn{2}{c}{Braço [m]} \\
\cmidrule(lr){6-7}\cmidrule(lr){8-9}
Parcela & Área [m$^2$] & $z$ [m] & $C_s$ & $C_h$ & 51,5 m/s & 36 m/s & Transv. & Long. \\
\midrule
Colunas (4) & 1.297,3 & 12,92 & 1,0 & 1,00 & 2.102,3 & 1.027,3 & 29,01 & 25,10 \\
\emph{Deck box} & 894,4 & 29,84 & 1,0 & 1,10 & 1.594,3 & 779,1 & 45,93 & 42,02 \\
Módulos & 1.341,6 & 39,84 & 1,1 & 1,20 & 2.869,8 & 1.402,3 & 55,93 & 52,02 \\
Face inferior & 12.499,2 & 25,84 & 1,0 & 1,10 & 22.280,8 & 10.887,3 & 41,93 & 38,02 \\
\bottomrule
\end{tabular}
\end{table}

\begin{table}[!htb]
\centering
\caption{Momentos de vento por parcela na solução final [kN$\cdot$m].}\label{tab:momentos}
\small
\begin{tabular}{lrrrr}
\toprule
 & \multicolumn{2}{c}{Transversal} & \multicolumn{2}{c}{Longitudinal} \\
\cmidrule(lr){2-3}\cmidrule(lr){4-5}
Parcela & 51,5 m/s & 36 m/s & 51,5 m/s & 36 m/s \\
\midrule
Colunas & 60.984 & 29.799 & 52.772 & 25.787 \\
\emph{Deck box} & 73.223 & 35.780 & 66.995 & 32.737 \\
Módulos & 160.499 & 78.426 & 149.290 & 72.949 \\
\midrule
$HM_0$ & 294.706 & 144.005 & 269.057 & 131.473 \\
$HM_{ud}$ (face inferior) & 934.163 & 456.471 & 847.134 & 413.945 \\
\bottomrule
\end{tabular}
\end{table}

\begin{table}[!htb]
\centering
\caption{Critério de áreas da solução final nos quatro casos.}\label{tab:criterio}
\small\setlength{\tabcolsep}{4pt}
\begin{tabular}{llrrrrrrrc}
\toprule
Plano & Vento & $z_{CLR}$ [m] & $HA(0)$ [m] & $\varphi_1$ [$^\circ$] & $\varphi_2$ [$^\circ$] &
$\phial$ [$^\circ$] & $A$ [m$\cdot$rad] & $B$ [m$\cdot$rad] & $A/B$ \\
\midrule
Transv. & 51,5 m/s & 7,89 & 0,746 & 14,12 & 31,62 & 21,57 & 0,574 & 0,440 & \textbf{1,303} \\
Transv. & 36 m/s & 7,89 & 0,365 & 6,68 & 33,12 & 21,57 & 0,574 & 0,215 & 2,667 \\
Long. & 51,5 m/s & 11,80 & 0,682 & 11,50 & 31,66 & 23,86 & 0,598 & 0,459 & 1,304 \\
Long. & 36 m/s & 11,80 & 0,333 & 5,04 & 33,03 & 23,86 & 0,598 & 0,224 & 2,669 \\
\bottomrule
\end{tabular}
\end{table}

Nos quatro casos, $\varphi_2 > \phial$ e a GZ é positiva de $0^\circ$ a $\phial$, e o critério é
atendido. O caso crítico é a tempestade no plano transversal, com $A/B = 1{,}303$, mas o
longitudinal fica praticamente no mesmo valor (1,304), sinal de que a otimização equilibrou os dois
planos. O plano transversal é o mais exigido porque a silhueta lateral dos \emph{pontoons}
($L_p h_p$) baixa o $z_{CLR}$ (7,89 contra 11,80~m) e alonga os braços.

As tabelas também mostram de onde vem o momento de vento. A face inferior do convés gera um momento
3,2 vezes maior que o das parcelas laterais, e como essa área é $L_p^2$, aumentar o convés aumenta o
emborcamento; é isso que mantém $L_p$ perto do mínimo. Entre as parcelas laterais, os módulos
respondem por 54\% de $HM_0$, o que explica a sensibilidade a $h_m$ (Seção~\ref{sec:sens-hdhm}).

A forma da curva GZ completa o quadro. O aumento rápido depois de cerca de 15$^\circ$ ocorre quando
o \emph{pontoon} do lado alto sai da água, e o centro de carena passa a acompanhar o \emph{pontoon}
imerso; depois de $\phial$ a curva não tem significado físico. Sob vento de tempestade, a plataforma
inclina 14,1$^\circ$ em equilíbrio no plano transversal, e 6,7$^\circ$ sob vento de operação. O
primeiro valor é elevado e merece atenção nas fases seguintes do projeto.

\subsection{\emph{Heave}}

A Figura~\ref{fig:forcas} decompõe a força de excitação da solução final nas três parcelas da
Seção~\ref{sec:heave-teoria}, e a Figura~\ref{fig:heave} mostra o RAO e o espectro de resposta. Os
valores principais estão na Tabela~\ref{tab:heave}.

\begin{figure}[!htb]
\centering
\includegraphics[width=0.7\textwidth]{\RBA/f_forcas_heave.png}
\caption[Solução final: parcelas da força de excitação em \emph{heave}]{Solução final: parcelas da
força de excitação em \emph{heave} por amplitude de onda, já multiplicadas pelo fator de
defasagem.}\label{fig:forcas}
\end{figure}

\begin{figure}[!htb]
\centering
\includegraphics[width=0.49\textwidth]{\RBA/f_rao.png}\hfill
\includegraphics[width=0.49\textwidth]{\RBA/f_espectros.png}
\caption[Solução final: RAO e espectro de resposta em \emph{heave}]{Solução final: RAO de
\emph{heave} (esquerda) e espectro de resposta com o espectro do mar de 100 anos
(direita).}\label{fig:heave}
\end{figure}

\begin{table}[!htb]
\centering
\caption{Resultados de \emph{heave} da solução final (caso-base).}\label{tab:heave}
\begin{tabular}{lr@{\hspace{1.2cm}}lr}
\toprule
Grandeza & Valor & Grandeza & Valor \\
\midrule
$\Delta$ [t] & 40.247 & $\omega_n$ [rad/s] & 0,2876 \\
$C_A$ & 1,451 & $T_n$ [s] & 21,85 \\
$l_d$ [m] & 86,69 & $\omega$ de cancelamento [rad/s] & 0,314 \\
$A_{33}$ [t] & 36.378 & RAO máximo & 1,47 \\
$C_{33}$ [kN/m] & 6.337 & RAO no pico do mar & 0,32 \\
$B_{33}$ [kN$\cdot$s/m] & 2.204 & $m_{0z}$ [m$^2$] & 0,1506 \\
 & & $\Hz$ ($\gamma = 3{,}3$ / $\gamma = 1$) [m] & 1,552 / 1,454 \\
\bottomrule
\end{tabular}
\end{table}

O período natural de 21,85~s fica 9\% acima do limite de 20~s. A massa adicional, de 36.378~t ou
90\% do deslocamento, quase dobra a massa efetiva e alonga o período natural em 38\% em relação ao
valor sem ela. O RAO tem pico de 1,47 em $\omega_n = 0{,}288$~rad/s, mas a frequência de
cancelamento, $\omega_c = 0{,}314$~rad/s, em que a parcela das colunas equilibra as dos
\emph{pontoons}, fica logo acima de $\omega_n$ e limita a ressonância. No pico do mar
($\omega_p = 0{,}413$~rad/s), o fator de defasagem vale 0,654, as parcelas valem
$F_{col} = 3.257$, $F_{pont} = -2.526$ e $F_{dif} = -2.928$~kN/m, e o RAO é de 0,32. Por isso o
espectro de resposta tem um pico pequeno em $\omega_n$ e o pico principal em $\omega_p$, e o
$\Hz$ resulta em 1,552~m, 28\% de $H_s$.

\subsection{\emph{Air gap}}

A Tabela~\ref{tab:airgap} aplica as expressões da Seção~\ref{sec:airgap-teoria}. A folga de 19,2~m é
grande porque o mar do campo é brando: a altura das colunas é definida pelo ângulo de alagamento
exigido pelo critério de vento, e não pela crista da onda.

\begin{table}[!htb]
\centering
\caption{Verificação do \emph{air gap} da solução final (caso-base).}\label{tab:airgap}
\begin{tabular}{lr}
\toprule
Grandeza & Valor [m] \\
\midrule
$\eta_{max} = 1{,}86 \times 5{,}6/2$ & 5,208 \\
$z_{min} = -1{,}86 \times 1{,}552/2$ & $-1{,}444$ \\
$\amin = \eta_{max} - z_{min}$ & 6,652 \\
$a = h_p + h_c - T$ & 25,837 \\
Folga $a - \amin$ & 19,186 \\
\bottomrule
\end{tabular}
\end{table}

\subsection{Quadro das restrições}

A Tabela~\ref{tab:resumo} reúne o valor de cada restrição na solução final e a folga em relação ao
limite. Ficam ativas exatamente as restrições que governam toda a fronteira, o $GM_T$ máximo e o
critério de áreas, além do lastro positivo. As tabelas equivalentes dos cenários 1 e 2 estão em
\path{caso_base/resultados/restricoes_1.csv} e \path{restricoes_2.csv}.

\begin{table}[!htb]
\centering
\caption[Restrições da solução final e folgas (caso-base)]{Valor de cada restrição na solução final
e folga em relação ao limite (caso-base).}\label{tab:resumo}
\small
\begin{tabular}{llrcrrrl}
\toprule
Restrição & Unid. & Valor & & Limite & Folga & Folga [\%] & \\
\midrule
Área de convés $L_p^2$ & m$^2$ & 12.499,2 & $\ge$ & 11.562,8 & 936,4 & 8,1 & \\
$GM_T$ mínimo & m & 4,000 & $\ge$ & 1,000 & 3,000 & 300,0 & \\
$GM_T$ máximo & m & 4,000 & $\le$ & 4,000 & 0,000 & 0,0 & \textbf{ativa} \\
$T_n$ & s & 21,85 & $>$ & 20,00 & 1,85 & 9,2 & \\
Lastro positivo & t & 11,0 & $>$ & 0 & 11,0 & 0,1 & \textbf{ativa} \\
$V_{lastro} \le 0{,}9 V_{pontoons}$ & m$^3$ & 10,7 & $\le$ & 27.524,5 & 27.513,7 & 99,96 & \\
$b_c \le b_p$ & m & 12,55 & $\le$ & 13,40 & 0,85 & 6,3 & \\
\emph{Air gap} $a \ge \amin$ & m & 25,84 & $\ge$ & 6,65 & 19,19 & 288,4 & \\
$T - h_p$ (geométrica) & m & 13,78 & $\ge$ & 1,00 & 12,78 & 1.277,6 & \\
$\min A/B$ (4 casos) & -- & 1,303 & $\ge$ & 1,300 & 0,003 & 0,2 & \textbf{ativa} \\
$GZ > 0$: $\min(GZ/\sin\varphi)$ & m & 4,000 & $>$ & 0 & 4,000 & -- & \\
\bottomrule
\end{tabular}
\end{table}

\begin{sintese}
A solução final tem casco de 14.998~t, $\Hz = 1{,}552$~m, $T_n = 21{,}85$~s, $GM_T = 4{,}00$~m e
$A/B = 1{,}303$, com folga de \emph{air gap} de 19,2~m. Ela atende a todas as restrições, e as ativas
são as de estabilidade, o que confirma que o projeto é governado pelo critério de vento e pelo
limite de GM.
\end{sintese}

% =============================================================================
\section{Sensibilidade às hipóteses}\label{sec:sensibilidade}

Algumas entradas do modelo são estimativas: as alturas do convés ($h_d$, $h_m$), a inclusão do gás
no BOE, a falta de lastro e os parâmetros hidrodinâmicos ($\gamma$, $\xi$). Esta seção mede quanto
cada uma muda o resultado, primeiro com a geometria da solução final fixa e depois repetindo a
otimização completa.

\subsection{Alturas do \emph{deck box} e dos módulos com geometria fixa}\label{sec:sens-hdhm}

Mantida a geometria da solução final, $h_d$ e $h_m$ foram variados em $h_d \in \{6, 8, 10\}$~m e
$h_m \in \{10, 12, 15\}$~m (Figura~\ref{fig:sensdh} e Tabela~\ref{tab:sensdh}).

\begin{figure}[!htb]
\centering
\includegraphics[width=0.9\textwidth]{\RBA/f_sens_hdhm.png}
\caption[Sensibilidade de $A/B$ e $GM_T$ a $h_d$ e $h_m$]{$\min A/B$ e $GM_T$ da solução final com
$h_d$ e $h_m$ alterados (geometria fixa). A estrela marca os valores adotados.}\label{fig:sensdh}
\end{figure}

\begin{table}[!htb]
\centering
\caption[KG, $GM_T$ e $A/B$ para outras alturas de convés]{KG, $GM_T$ e $\min A/B$ da solução final
para outras alturas de convés (geometria fixa).}\label{tab:sensdh}
\small
\begin{tabular}{rrrrrl}
\toprule
$h_d$ [m] & $h_m$ [m] & KG [m] & $GM_T$ [m] & $\min A/B$ & Viável? \\
\midrule
6 & 10 & 41,61 & 5,21 & 1,687 & não ($GM_T > 4$) \\
6 & 12 & 41,93 & 4,88 & 1,525 & não ($GM_T > 4$) \\
6 & 15 & 42,43 & 4,39 & 1,313 & não ($GM_T > 4$) \\
8 & 10 & 42,49 & 4,33 & 1,444 & não ($GM_T > 4$) \\
\textbf{8} & \textbf{12} & \textbf{42,82} & \textbf{4,00} & \textbf{1,303} & \textbf{sim} \\
8 & 15 & 43,31 & 3,51 & 1,107 & não ($A/B < 1{,}3$) \\
10 & 10 & 43,38 & 3,44 & 1,190 & não ($A/B < 1{,}3$) \\
10 & 12 & 43,70 & 3,12 & 1,063 & não ($A/B < 1{,}3$) \\
10 & 15 & 44,19 & 2,62 & 0,895 & não ($A/B < 1{,}3$) \\
\bottomrule
\end{tabular}
\end{table}

Só a combinação adotada ($h_d = 8$~m, $h_m = 12$~m) é viável sem redimensionar. Um convés mais
baixo reduz o KG e leva o $GM_T$ acima de 4~m; um convés mais alto aumenta o KG e a área exposta e
derruba o $A/B$. Esse comportamento é consequência direta das duas restrições ativas: a solução
está ajustada a elas, e qualquer mudança nas alturas exige reotimizar.

\subsection{Reotimizações}

Por isso, a otimização completa foi repetida para três hipóteses alternativas
(Tabela~\ref{tab:sensopt}), mudando uma variável de ambiente do código e a pasta de saída:
$h_m = 15$~m, valor do Exercício 9; gás incluído no BOE; e lastro mínimo de 10\% do volume dos
\emph{pontoons}.

\begin{table}[!htb]
\centering
\caption{Solução intermediária (joelho) de cada reotimização.}\label{tab:sensopt}
\small
\begin{tabular}{lrrrr}
\toprule
 & Caso-base & $h_m = 15$ m & Com gás & Lastro $\ge$ 10\% \\
\midrule
$L_p$ / $T$ [m] & 111,8 / 24,0 & 113,1 / 23,6 & 118,2 / 27,1 & 107,5 / 25,9 \\
$b_c$ / $h_c$ [m] & 12,55 / 39,6 & 12,83 / 41,7 & 13,66 / 42,0 & 12,85 / 36,5 \\
Peso do casco [t] & 14.998 & 15.792 & 18.718 & 16.026 \\
Variação & -- & +5,3\% & +24,8\% & +6,9\% \\
$\Hz$ [m] & 1,552 & 1,560 & 1,373 & 1,549 \\
$T_n$ [s] & 21,85 & 21,94 & 22,15 & 21,73 \\
Lastro [t] & 11 & 11 & 11 & 3.611 \\
Mínimo peso da fronteira [t] & 10.346 & 10.429 & 12.971 & 11.892 \\
\bottomrule
\end{tabular}
\end{table}

Com $h_m = 15$~m, o casco da solução escolhida fica 5,3\% mais pesado, com o mesmo $\Hz$, e as
conclusões qualitativas não mudam. Com o gás no BOE, o pico passa a 202,9 mil boe/dia, a planta a
16.838~t e a área mínima a 13.421~m$^2$ ($L_p \ge 115{,}85$~m); o casco da solução escolhida fica
24,8\% mais pesado (18.718~t), e o $\Hz$ cai porque a solução é maior e mais funda
($T = 27{,}1$~m). É a hipótese de maior impacto do projeto (Seção~\ref{sec:disc-boe}). Por fim, exigir
lastro mínimo de 10\% coloca 3.611~t nos \emph{pontoons}, abaixa o KG, custa 6,9\% de aço e mantém
$\Hz = 1{,}55$~m; essa condição é detalhada na Seção~\ref{sec:lastro}.

\subsection{Parâmetros do \emph{heave}}

Os parâmetros hidrodinâmicos foram variados com a geometria fixa (Tabela~\ref{tab:sensheave}).

\begin{table}[!htb]
\centering
\caption{$\Hz$ da solução final para diferentes $\xi$ e $\gamma$ (geometria fixa).}\label{tab:sensheave}
\begin{tabular}{lrrrr}
\toprule
$\xi$ & $\Hz$, $\gamma = 1$ [m] & $\Hz$, $\gamma = 3{,}3$ [m] & RAO máximo & $\amin$, $\gamma = 3{,}3$ [m] \\
\midrule
0,03 & 1,614 & 1,657 & 2,22 & 6,75 \\
0,05 & 1,454 & 1,552 & 1,47 & 6,65 \\
0,08 & 1,350 & 1,479 & 1,12 & 6,58 \\
\bottomrule
\end{tabular}
\end{table}

O $\Hz$ varia 6,7\% entre $\gamma = 1$ e $\gamma = 3{,}3$ e 12\% entre $\xi = 0{,}03$ e 0,08. O
amortecimento afeta sobretudo o pico do RAO em $\omega_n$, que cai de 2,22 para 1,12, e pouco a
resposta no pico do mar. No $\amin$, o efeito é menor que 0,2~m, desprezível diante da folga de
19~m.

\subsection{Síntese}

A Tabela~\ref{tab:sensresumo} resume as reotimizações. A hipótese dominante é a inclusão do gás no
BOE. As alturas do convés não podem mudar sem reotimizar, mas, reotimizada, a plataforma com
$h_m = 15$~m custa só 5,3\% a mais de casco. As incertezas hidrodinâmicas alteram o $\Hz$ em 7\% a
12\% e o \emph{air gap} mínimo em menos de 0,2~m.

\begin{table}[!htb]
\centering
\caption[Resumo da sensibilidade às hipóteses]{Resumo da sensibilidade às hipóteses (solução
escolhida de cada caso).}\label{tab:sensresumo}
\small
\begin{tabular}{lrrrr}
\toprule
Hipótese & Peso do casco & Variação & $\Hz$ & Lastro \\
\midrule
Caso-base (só óleo, $h_m = 12$~m) & 14.998 t & -- & 1,552 m & 11 t \\
$h_m = 15$~m (valor do Ex. 9) & 15.792 t & +5,3\% & 1,560 m & 11 t \\
Gás no BOE (202,9 mil boe/d) & 18.718 t & +24,8\% & 1,373 m & 11 t \\
Lastro mínimo de 10\% & 16.026 t & +6,9\% & 1,549 m & 3.611 t \\
\bottomrule
\end{tabular}
\end{table}


% =============================================================================
\section{Condição de 10\% de lastro}\label{sec:lastro}

A solução do caso-base tem um ponto fraco já apontado na Seção~\ref{sec:escolha}: o lastro é de só
11~t, uma margem numérica. Esta seção repete o projeto inteiro com lastro mínimo obrigatório e
compara o resultado com o caso-base. As tabelas marcadas como ``gerada'' vêm diretamente dos CSVs
das duas versões, pelo script \path{analises_condicao_10pct_lastro.py}.

\subsection{Motivação e formulação}

No caso-base, casco, \emph{topside}, \emph{risers} e amarração consomem todo o deslocamento. Uma
plataforma assim não teria como ajustar o calado entre trânsito e operação, corrigir banda e trim
causados por cargas variáveis nem compensar o crescimento de peso ao longo da vida útil. A
condição de 10\% de lastro acrescenta, por isso, uma restrição obrigatória,
\begin{equation}
V_{lastro} \ge 0{,}10\,V_{pontoons}, \qquad V_{pontoons} = 2 L_p b_p h_p ,
\end{equation}
que exige pelo menos 10\% do volume interno dos \emph{pontoons} ocupado por lastro na condição de
projeto. A fração foi aplicada ao volume dos \emph{pontoons}, onde o lastro fica, e não ao
deslocamento. Todo o resto (dados, modelo, validação, domínio, objetivos e demais restrições)
permanece igual ao caso-base.

O efeito da restrição não é óbvio de antemão. Por um lado, o lastro precisa de deslocamento, porque
$\Delta = \sum W$, e o casco tem de crescer. Por outro, a água fica no fundo dos \emph{pontoons},
com centro em $n/2$, e baixa o KG. Com KG menor, o $GM_T = KB + BM_T - KG$ chega a 4~m com um BM
menor, e a GZ cresce mais em ângulos grandes, o que facilita o critério de áreas. O saldo em peso de
casco só aparece com a otimização.

No código, a mudança se resume a um valor padrão: \texttt{LASTRO\_MIN\_FRAC = 0.10}, que acrescenta
a restrição em \texttt{tabela\_restricoes()}. As outras diferenças trocam a pasta de saída e os
rótulos, como mostra o \emph{diff} completo da \lst{diff}.

\subsection{Variação sistemática e fronteira de Pareto}

O procedimento é o mesmo da Seção~\ref{sec:procedimento}. Das 20.000 amostras do LHS, com a mesma
semente, 517 atendem às restrições estáticas, de \emph{heave}, de \emph{air gap} e de lastro, e 209
(1,0\%) também ao critério de áreas; o vento reprova sozinho 308 cascos, e a nova restrição tira 30
cascos do conjunto pré-viável. Na soma ponderada, os 21 valores de $w_1$ terminam em pontos viáveis,
mas 7 ficam dominados ($w_1$ = 0,05; 0,15; 0,45; 0,50; 0,55; 0,70 e 0,75), porque a fronteira
continua não convexa. Os 12 níveis do $\varepsilon$-restrito terminam em pontos viáveis com
\texttt{success = True}, e os dois mais altos convergem para o mesmo ponto. Das 242 soluções viáveis
(209 + 21 + 12), 23 são não dominadas, 12 da soma ponderada e 11 do $\varepsilon$-restrito
(Figura~\ref{fig:paretoL}; lista no Apêndice~\ref{ap:fronteiraL}). A rodada completa levou 176~s
em 8 processos.

\begin{figure}[!htb]
\centering
\includegraphics[width=0.88\textwidth]{\RL/01_pareto.png}
\caption{Conjunto viável e fronteira de Pareto na condição de 10\% de lastro.}\label{fig:paretoL}
\end{figure}

Como no caso-base, $GM_T = 4$~m e $A/B = 1{,}30$ estão ativos em todos os pontos
(Tabela~\ref{tab:ativasL}). A restrição de lastro está ativa em 16 dos 23: o otimizador põe
exatamente o lastro mínimo, e só os cascos mais profundos têm lastro acima dele. A área de convés
fica ativa em 6 pontos, contra 4 no caso-base, um primeiro sinal de que o lastro empurra $L_p$ para
o mínimo.

\begin{table}[!htb]
\centering
\caption[Restrições ativas ao longo da fronteira (10\% de lastro)]{Restrições ativas ao longo da
fronteira na condição de 10\% de lastro (gerada).}\label{tab:ativasL}
\begin{tabular}{lr}
\toprule
Restrição ativa & Pontos \\
\midrule
$GM_T = 4$~m & 23 / 23 \\
$A/B = 1{,}30$ & 23 / 23 \\
$V_{lastro} = 0{,}10\,V_{pontoons}$ & 16 / 23 \\
$b_c = b_p$ & 7 / 23 \\
área de convés mínima & 6 / 23 \\
\emph{Air gap} e $T_n$ & 0 / 23 \\
\bottomrule
\end{tabular}

\end{table}

A Figura~\ref{fig:variaveisL} mostra que as duas famílias de projeto do caso-base se mantêm. Até
15.995~t, a fronteira é formada por cascos rasos e largos, com $T$ de 12 a 22~m, $L_p$ crescendo até
150~m e $T_n$ de 31 a 46~s. A partir de 16.026~t, passa a cascos compactos e profundos, com $L_p$ de
107,5 a 116~m, $T$ de 26 a 40~m e $T_n$ de 21,7 a 23~s. A solução escolhida é o primeiro ponto da
família compacta.

\begin{figure}[!htb]
\centering
\includegraphics[width=0.85\textwidth]{\RLC/12_variaveis_fronteira.png}
\caption[Variáveis ao longo da fronteira (10\% de lastro)]{Variáveis normalizadas pelos limites ao
longo da fronteira (condição de 10\% de lastro).}\label{fig:variaveisL}
\end{figure}

\subsection{Três cenários}

A Tabela~\ref{tab:fichaL} e a Figura~\ref{fig:geo3L} apresentam os três cenários desta condição. O
de mínimo peso, com 11.892~t, continua sendo uma semissubmersível rasa, com $b_c$, $h_c$ e $T$ no
limite mínimo, mas agora com $L_p = 141{,}6$~m e \emph{pontoons} de 19,4~m de largura para carregar
3.991~t de lastro. Tem $T_n = 46{,}5$~s, mas o maior \emph{heave} dos três ($\Hz/H_s = 42\%$) e
ângulo de alagamento de só 8,7$^\circ$ (Figura~\ref{fig:gz1L}). O de mínimo \emph{heave}, com
23.634~t, é o mesmo casco do caso-base: com $T = 40$~m e $h_c = 50$~m, ele já tinha 18.081~t de
lastro (37,5\% dos \emph{pontoons}), e a nova restrição não o afeta. A solução escolhida, com
16.026~t, é detalhada a seguir.

\needspace{0.75\textheight}
{\small\renewcommand{\arraystretch}{1.04}
\begin{longtable}{lrrr}
\caption[Ficha técnica dos três cenários (10\% de lastro)]{Ficha técnica dos três cenários na
condição de 10\% de lastro (gerada).}\label{tab:fichaL}\\
\toprule
 & 1. Mín. peso & 2. Mín. \emph{heave} & 3. Escolhida \\
\midrule
\endfirsthead
\toprule
 & 1. Mín. peso & 2. Mín. \emph{heave} & 3. Escolhida \\
\midrule
\endhead
Origem & $w_1 = 0$ & $w_1 = 1$ & $\varepsilon = 1{,}549$ m \\
\linhas{tabelas_lastro10/tab_ficha_linhas.tex}
\bottomrule
\end{longtable}}

\begin{figure}[!htb]
\centering
\includegraphics[width=\textwidth]{\RL/05_geometria_3_cenarios.png}
\caption{Vista lateral dos três cenários na condição de 10\% de lastro.}\label{fig:geo3L}
\end{figure}

\begin{figure}[!htb]
\centering
\includegraphics[width=\textwidth]{\RL/04a_gz_min_peso.png}
\caption{Cenário 1 na condição de 10\% de lastro: curvas GZ e HA.}\label{fig:gz1L}
\end{figure}

\subsection{Solução final com lastro mínimo de 10\%}

\paragraph{Dimensões, pesos e centros}
A solução escolhida tem $L_p = 107{,}53$~m, no mínimo dado pela área de convés, $b_p = 13{,}17$~m,
$h_p = 12{,}43$~m, $b_c = 12{,}85$~m, $h_c = 36{,}52$~m e $T = 25{,}90$~m
(Tabela~\ref{tab:pesosL}). As 3.611~t de lastro ocupam 10\% do volume dos \emph{pontoons}, com nível
de 1,25~m e centro a 0,62~m da quilha. Com KB = 8,83~m e $KM_T = 42{,}38$~m, resultam
$GM_T = 4{,}00$~m e $GM_L = 4{,}22$~m, e o KG, de 38,38~m, fica 4,44~m abaixo do caso-base.

\begin{table}[!htb]
\centering
\caption[Pesos e centros da solução final (10\% de lastro)]{Pesos e centros da solução final na
condição de 10\% de lastro (gerada).}\label{tab:pesosL}
\begin{tabular}{lrrr}
\toprule
Item & Peso [t] & VCG [m] & Momento [t$\cdot$m] \\
\midrule
Planta de produção & 13.188,2 & 62,96 & 830.293 \\
Estrutura do convés & 4.615,9 & 52,96 & 244.444 \\
Carga operacional & 4.615,9 & 52,96 & 244.444 \\
\emph{Pontoons} & 9.509,1 & 6,22 & 59.118 \\
Colunas & 6.516,9 & 30,70 & 200.040 \\
\emph{Risers} & 1.808,9 & 48,96 & 88.559 \\
Amarração & 1.356,7 & 48,96 & 66.419 \\
Lastro (nos pontoons) & 3.611,0 & 0,62 & 2.246 \\
\midrule
\textbf{Total ($=\Delta$)} & \textbf{45.222,6} & \textbf{38,38} & \textbf{1.735.562} \\
\bottomrule
\end{tabular}

\end{table}

\paragraph{Estabilidade e critério de vento}
O critério é atendido nos quatro casos (Figura~\ref{fig:gzfinalL} e Tabelas~\ref{tab:ventofL}
a~\ref{tab:criterioL}), com mínimo de 1,303 sob tempestade no plano transversal e 1,304 no
longitudinal. Em todos os casos $\varphi_2 > \phial$ e a GZ é positiva até $\phial$, que vale
20,7$^\circ$ no plano transversal e 22,7$^\circ$ no longitudinal. A face inferior do convés gera 3,0
vezes o momento das parcelas laterais. Sob tempestade, a plataforma inclina 12,0$^\circ$ em
equilíbrio no plano transversal, contra 14,1$^\circ$ no caso-base, porque o convés menor e o
deslocamento maior reduzem o braço de vento inicial, $HA(0)$, de 0,746 para 0,612~m.

\begin{figure}[!htb]
\centering
\includegraphics[width=\textwidth]{\RL/02_gz_ha_solucao_final.png}
\caption[Solução final com 10\% de lastro: curvas GZ e HA]{Solução final com lastro mínimo de 10\%:
curvas GZ e HA nos dois planos ($\bullet$ $\varphi_1$, $\blacksquare$ $\varphi_2$).}\label{fig:gzfinalL}
\end{figure}

\begin{table}[!htb]
\centering
\caption[Forças e braços de vento (10\% de lastro)]{Forças e braços de vento por parcela, condição
de 10\% de lastro (gerada).}\label{tab:ventofL}
\small\setlength{\tabcolsep}{4.5pt}
\begin{tabular}{lrrrrrrrr}
\toprule
 & & & & & \multicolumn{2}{c}{Força [kN]} & \multicolumn{2}{c}{Braço [m]} \\
\cmidrule(lr){6-7}\cmidrule(lr){8-9}
Parcela & Área [m$^2$] & $z$ [m] & $C_s$ & $C_h$ & 51,5 m/s & 36 m/s & Transv. & Long. \\
\midrule
Colunas (4) & 1.185,4 & 11,53 & 1,0 & 1,00 & 1.920,9 & 938,6 & 28,55 & 24,56 \\
\emph{Deck box} & 860,2 & 27,05 & 1,0 & 1,10 & 1.533,5 & 749,3 & 44,08 & 40,08 \\
Módulos & 1.290,4 & 37,05 & 1,1 & 1,20 & 2.760,2 & 1.348,8 & 54,08 & 50,08 \\
Face inferior & 11.562,9 & 23,05 & 1,0 & 1,10 & 20.611,7 & 10.071,7 & 40,08 & 36,08 \\
\bottomrule
\end{tabular}

\end{table}

\begin{table}[!htb]
\centering
\caption[Momentos de vento (10\% de lastro)]{Momentos de vento por parcela [kN$\cdot$m], condição de
10\% de lastro (gerada).}\label{tab:ventomL}
\small
\begin{tabular}{lrrrr}
\toprule
 & \multicolumn{2}{c}{Transversal} & \multicolumn{2}{c}{Longitudinal} \\
\cmidrule(lr){2-3}\cmidrule(lr){4-5}
Parcela & 51,5 m/s & 36 m/s & 51,5 m/s & 36 m/s \\
\midrule
Colunas (4) & 54.839 & 26.797 & 47.172 & 23.050 \\
\emph{Deck box} & 67.589 & 33.027 & 61.468 & 30.036 \\
Módulos & 149.262 & 72.936 & 138.245 & 67.552 \\
\midrule
$HM_0$ & 271.690 & 132.759 & 246.886 & 120.639 \\
$HM_{ud}$ (face inferior) & 826.039 & 403.637 & 743.769 & 363.437 \\
\bottomrule
\end{tabular}

\end{table}

\begin{table}[!htb]
\centering
\caption[Critério de áreas (10\% de lastro)]{Critério de áreas nos quatro casos, condição de 10\% de
lastro (gerada).}\label{tab:criterioL}
\small\setlength{\tabcolsep}{4pt}
\begin{tabular}{llrrrrrrrc}
\toprule
Plano & Vento & $z_{CLR}$ [m] & $HA(0)$ [m] & $\varphi_1$ [$^\circ$] & $\varphi_2$ [$^\circ$] &
$\phial$ [$^\circ$] & $A$ [m$\cdot$rad] & $B$ [m$\cdot$rad] & $A/B$ \\
\midrule
\linhas{tabelas_lastro10/tab_criterio_linhas.tex}
\bottomrule
\end{tabular}
\end{table}

\paragraph{\emph{Heave}}
O período natural de 21,73~s fica 8,6\% acima do limite de 20~s (Figura~\ref{fig:heaveL} e
Tabela~\ref{tab:heaveL}). A força de excitação se anula em 0,318~rad/s, logo acima de
$\omega_n = 0{,}289$~rad/s, o que limita o pico do RAO a 1,52. No pico do mar, o fator de
defasagem vale 0,680 e o RAO é de 0,32, e o espectro de resposta tem, como no caso-base, um pico
pequeno em $\omega_n$ e o principal no pico da onda. O $\Hz$ resulta em 1,549~m (28\% de $H_s$),
praticamente o do caso-base.

\begin{figure}[!htb]
\centering
\includegraphics[width=\textwidth]{\RL/03_heave_solucao_final.png}
\caption[Solução final com 10\% de lastro: \emph{heave}]{Solução final com lastro mínimo de 10\%:
força de excitação, RAO, espectro do mar e espectro de resposta.}\label{fig:heaveL}
\end{figure}

\begin{table}[!htb]
\centering
\caption[Resultados de \emph{heave} (10\% de lastro)]{Resultados de \emph{heave}, condição de 10\% de
lastro (gerada).}\label{tab:heaveL}
\begin{tabular}{lr@{\hspace{1.5cm}}lr}
\toprule
Grandeza & Valor & Grandeza & Valor \\
\midrule
$\Delta$ [t] & 45.223 & $\omega_n$ [rad/s] & 0,2892 \\
$C_A$ & 1,498 & $T_n$ [s] & 21,73 \\
$l_d$ [m] & 81,82 & $\omega$ de cancelamento [rad/s] & 0,318 \\
$A_{33}$ [t] & 34.223 & RAO máximo & 1,52 \\
$C_{33}$ [kN/m] & 6.645 & RAO no pico do mar & 0,32 \\
$B_{33}$ [kN$\cdot$s/m] & 2.298 & $m_{0z}$ [m$^2$] & 0,1500 \\
 & & $\Hz$ ($\gamma = 3{,}3$ / $\gamma = 1$) [m] & 1,549 / 1,456 \\
\bottomrule
\end{tabular}

\end{table}

\paragraph{\emph{Air gap} e restrições}
A folga de \emph{air gap} cai para 16,4~m, ainda ampla (Tabela~\ref{tab:airgapL}). No quadro das
restrições (Tabela~\ref{tab:resumoL}), ficam ativas a área de convés, com $L_p$ no mínimo, o $GM_T$
máximo, o lastro mínimo de 10\% e o critério de áreas.

\begin{table}[!htb]
\centering
\caption[\emph{Air gap} (10\% de lastro)]{Verificação do \emph{air gap}, condição de 10\% de lastro
(gerada).}\label{tab:airgapL}
\begin{tabular}{lr}
\toprule
Grandeza & Valor [m] \\
\midrule
$\eta_{max} = 1{,}86 \times 5{,}6/2$ & 5,208 \\
$z_{min} = -1{,}86 \times 1,549/2$ & $-1,441$ \\
$a_{\text{mín}} = \eta_{max} - z_{min}$ & 6,649 \\
$a = h_p + h_c - T$ & 23,055 \\
Folga $a - a_{\text{mín}}$ & 16,406 \\
\bottomrule
\end{tabular}

\end{table}

\begin{table}[!htb]
\centering
\caption[Restrições da solução final (10\% de lastro)]{Valor de cada restrição na solução final e
folga, condição de 10\% de lastro (gerada).}\label{tab:resumoL}
\small
\begin{tabular}{llrcrrrl}
\toprule
Restrição & Unid. & Valor & & Limite & Folga & Folga [\%] & \\
\midrule
\linhas{tabelas_lastro10/tab_restricoes_linhas.tex}
\bottomrule
\end{tabular}
\end{table}

\paragraph{Escolha}
A solução escolhida também é o joelho desta fronteira. Em relação ao cenário de mínimo peso, aumenta
o casco em 34,8\% e reduz o $\Hz$ em 34,3\%; em relação ao de mínimo \emph{heave}, economiza 32,2\%
do aço com um $\Hz$ 48,5\% maior. Entre os extremos, cortar o \emph{heave} em 56\% custa 99\% a mais
de peso, e a solução escolhida captura 62\% do ganho de movimento com 35\% do acréscimo de massa.
Depois dela, cada milímetro de \emph{heave} custa 2,9 vezes mais aço (Tabela~\ref{tab:custoL}), e o
ângulo de alagamento, de 20,7$^\circ$, é bem maior que os 8,7$^\circ$ do cenário de mínimo peso. Ao
contrário do caso-base, uma variável fica no limite, $L_p$ no mínimo, mas esse limite é a própria
restrição de área de convés, e não uma escolha de domínio.

\begin{table}[!htb]
\centering
\caption[Custo marginal de redução do \emph{heave} (10\% de lastro)]{Custo marginal de redução do
\emph{heave}, condição de 10\% de lastro (gerada).}\label{tab:custoL}
\begin{tabular}{lrrr}
\toprule
Trecho & Peso adicional & Redução de $\Hz$ & Custo [t/mm] \\
\midrule
\linhas{tabelas_lastro10/tab_custo_linhas.tex}
\bottomrule
\end{tabular}
\end{table}

\subsection{Comparação com o caso-base}\label{sec:lastro-comp}

As Tabelas~\ref{tab:cenbase} e~\ref{tab:comparacao} e a Figura~\ref{fig:fronteiras} colocam as duas
versões lado a lado.

\begin{table}[!htb]
\centering
\caption[Os três cenários nas duas versões]{Os três cenários no caso-base e na condição de 10\% de
lastro (gerada).}\label{tab:cenbase}
\begin{tabular}{lrrrr}
\toprule
Cenário & \multicolumn{2}{c}{Peso do casco [t]} & \multicolumn{2}{c}{$\Hz$ [m]} \\
\cmidrule(lr){2-3}\cmidrule(lr){4-5}
 & Caso-base & Lastro $\ge$ 10\% & Caso-base & Lastro $\ge$ 10\% \\
\midrule
1 · Mín. peso & 10.346 & 11.892 & 2,366 & 2,359 \\
2 · Mín. \emph{heave} & 23.634 & 23.634 & 1,043 & 1,043 \\
3 · Escolhida & 14.998 & 16.026 & 1,552 & 1,549 \\
\bottomrule
\end{tabular}

\end{table}

\begin{figure}[!htb]
\centering
\includegraphics[width=0.82\textwidth]{\RLC/11_fronteiras_base_x_lastro10.png}
\caption{Fronteiras de Pareto do caso-base e da condição de 10\% de lastro.}\label{fig:fronteiras}
\end{figure}

\begin{table}[!htb]
\centering
\caption[Solução escolhida: caso-base contra 10\% de lastro]{Solução escolhida: caso-base contra
condição de 10\% de lastro (gerada).}\label{tab:comparacao}
\begin{tabular}{lrrr}
\toprule
Grandeza & Caso-base & Lastro $\ge$ 10\% & Variação \\
\midrule
\linhas{tabelas_lastro10/tab_comparacao_linhas.tex}
\bottomrule
\end{tabular}
\end{table}

A primeira observação é que a comparação é feita a igual \emph{heave}. As duas soluções escolhidas
vêm do mesmo nível do $\varepsilon$-restrito, o quinto de doze ($\varepsilon = 1{,}5523$~m no
caso-base e 1,5494~m com lastro), porque os extremos de $\Hz$ das duas fronteiras são quase iguais
(2,366 e 2,359~m de um lado e 1,043~m dos dois lados do outro). Como a restrição
$\Hz \le \varepsilon$ está ativa nas duas, o $\Hz$ (1,552 contra 1,549~m) e o $T_n$ (21,85 contra
21,73~s) praticamente coincidem, e a diferença de peso mede o custo do lastro a igual
\emph{heave}.

Esse custo é menor do que o lastro faria supor. As 3.611~t de lastro exigem 4.976~t a mais de
deslocamento (+12,4\%): 3.600~t de lastro, 1.028~t de casco e 349~t de \emph{risers} e amarração. O
casco cresce só 6,9\% porque o volume extra vem de \emph{pontoons} mais altos ($h_p$ +2,2~m,
+1.252~t de aço; cada m$^3$ de \emph{pontoon} desloca 1,025~t e custa 0,27~t de aço), enquanto as
colunas ficam 224~t mais leves. A água no fundo baixa o KG em 4,44~m; com KG mais baixo, o
$GM_T = 4$~m é obtido com BM menor (33,6 contra 39,1~m), o que aproxima as colunas e leva $L_p$ ao
mínimo. A GZ também cresce mais em ângulos grandes, e a área $A$ sai com colunas mais curtas ($h_c$
$-$3,1~m), mantendo $GM_T$ e $A/B$ ativos. A folga de \emph{air gap}, por sua vez, cai 2,8~m.

Na fronteira como um todo, a restrição desloca a parte leve para a direita, com o mínimo peso
subindo 14,9\%, de 10.346 para 11.892~t, e deixa intacto o extremo de mínimo \emph{heave}, que já
tinha lastro. Perto da solução escolhida, a fronteira fica cerca de 1.000~t à direita da do
caso-base.

\subsection{Comparação com a frota real}

O lastro também aproxima a plataforma da frota real (Tabela~\ref{tab:posterL} e
Figura~\ref{fig:posterL}). O deslocamento específico sobe para 285~t por mil bopd, o nível da Troll C
(286~t), enquanto no caso-base era de 253~t, abaixo de toda a frota. Isso confirma que a leveza do
caso-base vinha da falta de lastro. Calado, altura e lado da coluna e largura do \emph{pontoon}
ficam dentro da faixa da frota; o \emph{pontoon}, de 12,4~m, fica mais alto que o maior do pôster
(10,1~m), efeito do volume extra para o lastro.

\begin{table}[!htb]
\centering
\caption[Solução com 10\% de lastro contra a frota do pôster]{Solução escolhida (10\% de lastro)
contra as semissubmersíveis do pôster de 2011 (gerada).}\label{tab:posterL}
\small
\begin{tabular}{lrrrrr}
\toprule
Grandeza & $n$ & Mínimo & Mediana & Máximo & G02 (10\%) \\
\midrule
\linhas{tabelas_lastro10/tab_poster_linhas.tex}
\bottomrule
\end{tabular}
\end{table}

\begin{figure}[!htb]
\centering
\includegraphics[width=\textwidth]{\RLC/13_poster2011.png}
\caption[Solução com 10\% de lastro contra o pôster de 2011]{Solução escolhida com lastro mínimo de
10\% contra o pôster de 2011: (a) dimensões do casco; (b) deslocamento por mil bopd.}\label{fig:posterL}
\end{figure}

\begin{sintese}
As duas configurações são tecnicamente válidas e equivalentes em hidrodinâmica ($\Hz = 1{,}55$~m e
$T_n \approx 21{,}8$~s). Incluir 3.611~t de lastro custa 6,9\% a mais de aço (16.026~t), reduz o KG
em 4,44~m e alinha o deslocamento à frota real. A condição de 10\% de lastro é a recomendada para a
fase seguinte do projeto.
\end{sintese}

% =============================================================================
\section{Discussão das hipóteses e limitações}\label{sec:discussao}

Os resultados das seções anteriores dependem de algumas entradas que o enunciado não fixa. Esta
seção discute cada uma delas e, em seguida, as limitações do próprio modelo.

\subsection{BOE só com óleo ou com gás}\label{sec:disc-boe}

O caso-base dimensiona a planta só pelo pico de óleo, de 158.894~bpd. Três argumentos sustentam
essa escolha: o enunciado trata de uma plataforma de produção de óleo, a curva de produção dada é
de óleo e, em reservatórios com capa de gás, o gás costuma ser reinjetado ou produzido depois do
óleo, de modo que os dois picos não coincidem. Em sentido contrário, unidades reais tratam o gás
associado a bordo, e somar o gás é o lado conservador para o peso. Incluído com o mesmo perfil de
produção, o gás eleva o pico a 202,9 mil boe/dia, a planta a 16.838~t e o casco da solução escolhida
em 24,8\% (Seção~\ref{sec:sensibilidade}). Como é a hipótese de maior impacto do projeto, ela deve
ser confirmada com o professor, assim como a rota de exportação do gás (gasoduto de 190~km ou
reinjeção), que o enunciado também não define.

\subsection{Peso da planta de 83~kg/boe}\label{sec:disc-planta}

O valor de 0,083~t por boe/dia vem do código original do grupo, e o grupo não identificou sua fonte
nos materiais da disciplina. Duas verificações indiretas, porém, o sustentam. Na plataforma de
referência dos Exercícios 7 a 10, com BOE = 120.000, ele dá uma planta de 9.960~t e um KG de
33,21~m, contra 33,00~m da aula; como o \emph{topside} domina o KG, um peso de planta muito
diferente do usado em aula apareceria nessa comparação. Além disso, o \emph{topside} resultante, de
141~t por mil bopd, coincide com a mediana das semissubmersíveis reais do pôster de 2011, de 145~t
por mil bopd.

\subsection{Limites das variáveis}\label{sec:disc-limites}

Só o limite inferior de $L_p$ decorre dos dados do campo. Os outros vêm do código original do grupo,
sem fonte no enunciado, e são mais amplos que a faixa da frota do pôster (Tabela~\ref{tab:limites}).
A consequência é que os extremos da fronteira dependem desses limites: o cenário 1 tem colunas de
8~m e calado de 12~m, e o cenário 2 tem calado de 40~m com colunas de 50~m, cascos sem precedente na
frota. A solução escolhida, ao contrário, tem todas as variáveis no interior do domínio e não depende
deles; ela também está dentro da faixa da frota, exceto pela altura do \emph{pontoon}, 0,1~m acima
do máximo do pôster. Uma alternativa seria usar a faixa do pôster como domínio, o que mudaria os
extremos da fronteira; essa rodada não foi feita.

\subsection{Alturas do \emph{deck box} e dos módulos}

O valor $h_d = 8$~m tem referência indireta no pôster (11~m para dois níveis na Independence Hub) e
foi mantido porque dá o melhor ajuste conjunto de KG e $A/B$ na validação. Já $h_m = 12$~m é uma
estimativa do grupo. A solução final só é viável com esses valores (Tabela~\ref{tab:sensdh}), mas a
reotimização com $h_m = 15$~m custa só 5,3\% de casco, e as conclusões qualitativas não dependem
dessas alturas.

\subsection{Limitações do modelo}\label{sec:disc-limitacoes}

O modelo paramétrico simplifica a plataforma em vários pontos, e alguns deles merecem atenção nas
fases seguintes:
\begin{itemize}
\item na estabilidade, a GZ 2D não inclui o volume do convés nem o acoplamento com o trim, e não há
verificação de estabilidade em avaria nem de superfície livre do lastro. A inclinação de equilíbrio
sob tempestade, de 14,1$^\circ$ no caso-base e 12,0$^\circ$ com lastro, atende ao critério, mas é
alta e deve ser verificada;
\item nos movimentos, só o \emph{heave} em ondas de través é calculado, \emph{roll} e \emph{pitch}
não entram no \emph{air gap}, e o amortecimento é uma fração fixa do crítico;
\item no \emph{air gap}, o cálculo é linear, sem a amplificação da crista pelas colunas nem cristas
não lineares (DNVGL-OTG-13~\cite{otg13}); com folgas de 16 a 19~m, essa limitação pouco importa no
mar do campo;
\item nos pesos, as correlações não têm margens de projeto para aço secundário, equipamentos de
casco e crescimento de peso;
\item no escopo, ancoragem, corrente, \emph{risers} e escoamento da produção (FSO ou oleoduto de
190~km) ficam fora do modelo paramétrico.
\end{itemize}

% =============================================================================
\section{Conclusões e recomendações}\label{sec:conclusoes}

O modelo paramétrico desenvolvido calcula todos os itens pedidos e é confiável para a escala do
projeto: reproduz os resultados das aulas com erros de 0,0\% a 3,7\% e o deslocamento de três
semissubmersíveis reais de 2 \emph{pontoons} e 4 colunas com erro de até 10\%.

Aplicado ao campo, ele mostra que o domínio de projeto é pouco viável, com só 1,1\% das 20.000
combinações amostradas atendendo a todas as restrições, e que o critério de áreas é a restrição mais
seletiva. A fronteira de Pareto do caso-base vai de 10.346~t ($\Hz = 2{,}37$~m) a 23.634~t
($\Hz = 1{,}04$~m), é não convexa e tem duas famílias de projeto, uma rasa e larga e outra compacta
e profunda; sem o $\varepsilon$-restrito, a soma ponderada deixaria vazios nessa fronteira. Em
todas as soluções ótimas, $GM_T = 4$~m e $A/B = 1{,}30$ estão ativos, enquanto \emph{air gap} e
$T_n$ têm folga: o projeto é governado pela estabilidade.

A solução escolhida pesa 14.998~t e tem $\Hz = 1{,}55$~m, $T_n = 21{,}85$~s, $A/B = 1{,}303$ e
$\phial = 21{,}6^\circ$. Ela está no ponto a partir do qual reduzir o \emph{heave} custa três vezes
mais aço e, no caso-base, tem todas as variáveis no interior do domínio. Suas dimensões e seu
\emph{topside} estão na faixa das semissubmersíveis reais, mas o deslocamento específico, de 253~t
por mil bopd, fica abaixo de toda a frota por falta de lastro. Na região, os análogos reais (CLOV e
Moho/Bilondo) têm casco de navio.

Com lastro mínimo de 10\% do volume dos \emph{pontoons}, a solução escolhida continua equivalente em
hidrodinâmica ($\Hz = 1{,}55$~m e $T_n \approx 21{,}8$~s). Incluir 3.611~t de lastro custa 6,9\% a
mais de aço (16.026~t), reduz o KG em 4,44~m e leva o deslocamento específico a 285~t por mil bopd,
o nível da Troll C. Entre as hipóteses testadas, a dominante é a inclusão do gás no BOE, que
aumentaria o casco em 24,8\%; as demais mudam o casco em 5\% a 7\% e o $\Hz$ em até 12\%.

Para a fase seguinte, o grupo recomenda:
\begin{itemize}
\item adotar a condição de 10\% de lastro, que garante ajuste de calado e de trim e margem para
crescimento de peso;
\item confirmar com o professor se o gás deve entrar no BOE;
\item fundamentar $h_m$ e o peso de planta de 83~kg/boe com dados de unidades semelhantes;
\item verificar a inclinação de 12 a 14$^\circ$ sob vento de tempestade, a estabilidade em avaria e
a superfície livre do lastro;
\item incluir \emph{roll}, \emph{pitch} e os efeitos não lineares no \emph{air gap} (DNVGL-OTG-13);
\item estender a validação com plataformas reais à estabilidade e ao \emph{heave}, que o pôster de
2011 não cobre;
\item incluir o custo do armazenamento ou do escoamento (FSO ou oleoduto de 190~km) na comparação
entre semissubmersível e FPSO, a solução usada na região.
\end{itemize}

% =============================================================================
\begin{thebibliography}{99}
\bibitem{chakrabarti} Chakrabarti, S. (Ed.). \emph{Handbook of Offshore Engineering}, Cap. 7,
Vol. 1. Elsevier, 2005.
\bibitem{goncalves} Gonçalves, R.T.; Matsumoto, F.T. \emph{Projeto de uma Unidade de Exploração
e Produção de Petróleo}. Trabalho final de PNV2613, USP, 2005.
\bibitem{brotzman} Brotzman, D. et al. \emph{Student Design of a Semi-submersible Floating
Production System for the Gulf of Mexico}. TSNAME, 2012.
\bibitem{gallala} Gallala, J.R. \emph{Hull Dimension of a Semisubmersible Rig -- A Parametric
Optimization Approach}. NTNU, 2013.
\bibitem{dsouza} D'Souza, R. et al. \emph{The Semisubmersible Floating Production System: A
25-Year Technical and Historic Retrospective}. SNAME Maritime Convention, 2019.
\bibitem{poster} Supan, C. (Mustang Engineering). \emph{2011 Worldwide Survey of Semi-FPSs and
FPUs}. Offshore Magazine, Poster 89, jan. 2011.
\bibitem{c301} DNV. \emph{DNV-OS-C301 -- Stability and Watertight Integrity}, 2013.
\bibitem{modu} IMO. \emph{Code for the Construction and Equipment of Mobile Offshore Drilling
Units (2009 MODU Code)}, Res. A.1023(26), 2009.
\bibitem{c205} DNV. \emph{DNV-RP-C205 -- Environmental Conditions and Environmental Loads}, 2010.
\bibitem{otg13} DNV GL. \emph{DNVGL-OTG-13 -- Prediction of Air Gap for Column-Stabilised
Units}, 2019.
\bibitem{mckay} McKay, M.D.; Beckman, R.J.; Conover, W.J. A comparison of three methods for
selecting values of input variables in the analysis of output from a computer code.
\emph{Technometrics}, 21(2), p. 239--245, 1979.
\bibitem{kraft} Kraft, D. \emph{A Software Package for Sequential Quadratic Programming}.
DFVLR-FB 88-28, Köln, 1988.
\bibitem{haimes} Haimes, Y.Y.; Lasdon, L.S.; Wismer, D.A. On a bicriterion formulation of the
problems of integrated system identification and system optimization. \emph{IEEE Transactions on
Systems, Man, and Cybernetics}, SMC-1(3), p. 296--297, 1971.
\bibitem{miettinen} Miettinen, K. \emph{Nonlinear Multiobjective Optimization}. Kluwer, 1999.
\bibitem{scipy} Virtanen, P. et al. SciPy 1.0: fundamental algorithms for scientific computing in
Python. \emph{Nature Methods}, 17, p. 261--272, 2020.
\bibitem{clov} Offshore Technology. \emph{CLOV Development Project, Block 17, Offshore Angola}.
\url{https://www.offshore-technology.com/projects/clov-development-project-angola/}; e TotalEnergies,
\emph{CLOV start-up} (2014).
\bibitem{scarabeo5} Saipem. \emph{Saipem completes conversion of Scarabeo 5 from drilling unit to
floating production facility for Eni in Congo}, 26/08/2025. Dimensões de 111 × 73~m conforme
registro AIS (VesselTracker, IMO 8754970).
\bibitem{omsemis} Beaubouef, B. \emph{Semisubmersibles setting the standard in the Gulf of
Mexico}. Offshore Magazine, dez. 2021.
\bibitem{omgom} Beaubouef, B. \emph{Cost efficiency underpins new Gulf of Mexico platform
designs}. Offshore Magazine, nov. 2019.
\end{thebibliography}


% =============================================================================
\appendix
\section{Implementação computacional}\label{ap:codigo}

Este apêndice complementa as seções de teoria com a implementação: a organização dos scripts, a
correspondência entre funções e equações, as medidas de desempenho e os trechos do código real.

\subsection{Estrutura do repositório e scripts}

O projeto foi implementado em Python~3 com NumPy, SciPy, pandas e Matplotlib, listados no arquivo
\path{requirements.txt}. A Tabela~\ref{tab:scripts} mostra os scripts e o que cada um produz.

\begin{table}[!htb]
\centering
\caption{Scripts do repositório e saídas.}\label{tab:scripts}
\small
\begin{tabular}{>{\raggedright\arraybackslash}p{6.2cm}>{\raggedright\arraybackslash}p{8.2cm}}
\toprule
Script & Função e saídas \\
\midrule
\path{caso_base/codigo/pnv3425_grupo02.py} & modelo, validação, variação sistemática, soma
ponderada, $\varepsilon$-restrito, fronteira, três cenários e solução final; figuras 00--05 e CSVs
em \path{caso_base/resultados/} \\
\path{caso_base/codigo/analises_complementares.py} & curva de produção, restrições ativas e
variáveis ao longo da fronteira, sensibilidade a $h_d$/$h_m$, $\xi$ e $\gamma$, custo marginal
(figuras 06--08) \\
\path{caso_base/codigo/comparativo_plataformas_reais.py} & CLOV, Scarabeo 5 e semis do Golfo do
México (figura 09) \\
\path{caso_base/codigo/comparativo_poster_2011.py} & validação do casco com o pôster de 2011 e
comparação com a frota (figura 10) \\
\path{caso_base/codigo/figuras_apresentacao.py} & figuras no padrão visual da apresentação
(\path{resultados/figuras_apresentacao/}) \\
\path{condicao_10pct_lastro/codigo/pnv3425_grupo02_condicao_10pct_lastro.py} & o mesmo modelo com
lastro mínimo de 10\%; saídas em \path{resultados_lastro10/} \\
\path{condicao_10pct_lastro/codigo/analises_condicao_10pct_lastro.py} & comparação com o caso-base
e com o pôster (figuras 11--13) e tabelas em \path{relatorio/tabelas_lastro10/} \\
\path{relatorio/gerar_trechos.py} & copia do código os trechos deste apêndice
(\path{relatorio/trechos/}) e gera o \emph{diff} entre as duas versões \\
\bottomrule
\end{tabular}
\end{table}

\subsection{Funções do modelo e equações}

A Tabela~\ref{tab:funcoes} liga cada função do modelo às grandezas que ela calcula, à seção em que
a teoria correspondente foi apresentada e à listagem com o código.

\begin{table}[!htb]
\centering
\caption{Funções do modelo, grandezas calculadas e seção da teoria.}\label{tab:funcoes}
\small
\begin{tabular}{>{\raggedright\arraybackslash}p{4.4cm}>{\raggedright\arraybackslash}p{6.9cm}rr}
\toprule
Função & Calcula & Seção & Listagem \\
\midrule
\texttt{Cenario}, \texttt{BOUNDS} & BOE, $Q$, área mínima, $\eta_{max}$; domínio &
\ref{sec:fixos}, \ref{sec:dominio} & \ref{lst:cenario} \\
\texttt{hidrostatica} & $\nabla$, $\Delta$, KB, $I_{T,L}$, $BM_{T,L}$, $KM_{T,L}$, $A_{wp}$ &
\ref{sec:hidro} & \ref{lst:hidrostatica} \\
\texttt{pesos} & itens de peso, $\Wc$, lastro e nível, KG, $GM_{T,L}$ & \ref{sec:pesos} &
\ref{lst:pesos} \\
\texttt{caixas\_2d}, \texttt{\_integrais\_submersas} & retângulos da projeção 2D; área e momentos
submersos & \ref{sec:gz-teoria} & \ref{lst:caixas} \\
\texttt{equilibrio\_2d}, \texttt{curva\_gz} & linha d'água a volume constante; $GZ(\varphi)$ &
\ref{sec:gz-teoria} & \ref{lst:equilibrio} \\
\texttt{angulo\_alagamento} & $\phial$ & \ref{sec:gz-teoria} & \ref{lst:alagamento} \\
\texttt{z\_clr}, \texttt{tabela\_vento} & $z_{CLR}$, forças, braços, $HM_0$, $HM_{ud}$ &
\ref{sec:vento-teoria} & \ref{lst:vento} \\
\texttt{criterio\_areas}, \texttt{estabilidade} & $HA$, $\varphi_1$, $\varphi_2$, $A$, $B$, $A/B$,
$\min(GZ/\sin\varphi)$ & \ref{sec:vento-teoria} & \ref{lst:criterio} \\
\texttt{ca\_retangulo}, \texttt{espectro\_mar}, \texttt{heave} & $C_A$, $A_{33}$, $C_{33}$, $B_{33}$,
$T_n$, $F_3$, RAO, $S_\zeta$, $S_z$, $\Hz$ & \ref{sec:heave-teoria}, \ref{sec:espectro-teoria} &
\ref{lst:heave} \\
\texttt{avaliar}, \texttt{tabela\_restricoes} & \emph{air gap} mínimo, restrições, folgas,
viabilidade & \ref{sec:airgap-teoria}, \ref{sec:restricoes} & \ref{lst:avaliar} \\
\texttt{variacao\_sistematica}, \texttt{nao\_dominados} & LHS em dois estágios; filtro de Pareto &
\ref{sec:procedimento} & \ref{lst:lhs} \\
\texttt{Problema}, \texttt{\_tarefa\_slsqp} & \emph{cache} e partidas do SLSQP & \ref{sec:procedimento}
& \ref{lst:problema} \\
\texttt{soma\_ponderada} & ótimos para $w_1 = 0; 0{,}05; \ldots; 1$ & \ref{sec:procedimento} &
\ref{lst:somapond} \\
\texttt{epsilon\_restrito}, \texttt{ponto\_joelho} & ótimos com $\Hz \le \varepsilon$; joelho &
\ref{sec:procedimento} & \ref{lst:epsilon} \\
\bottomrule
\end{tabular}
\end{table}

\subsection{Desempenho e robustez numérica}\label{sec:cod-desempenho}

A curva GZ é a parte mais cara do modelo, porque exige resolver a linha d'água a volume constante
em cada ângulo. Quatro medidas reduziram a rodada completa de 1.555~s para 329~s, e a rodada final
registrada no log levou 241~s em 8 processos. A primeira foi trocar a bissecção pura pelo método de
Newton com derivada exata, usando o calado $T$ como chute inicial. A segunda foi vetorizar o
cálculo, resolvendo todos os ângulos de uma curva ao mesmo tempo em operações de \emph{arrays}, com
ângulos nas linhas e retângulos nas colunas. A terceira foi tirar o \texttt{pandas} do laço: as
tabelas de pesos e de vento são listas durante a otimização e só viram \texttt{DataFrame} no
relatório. A quarta foi compartilhar as avaliações entre objetivo e restrições do SLSQP por meio de
um \emph{cache} e rodar as partidas de cada etapa em processos separados.

A troca para Newton trouxe, na primeira versão, picos isolados nas curvas GZ. A causa era o método
ciclar entre as duas pontas do intervalo de busca quando a linha d'água passava por um canto de
retângulo, onde $V(\alpha)$ tem quina e a derivada muda de valor. A versão final mantém, para cada
ângulo, o intervalo $[lo, hi]$ da bissecção, aceita o passo de Newton só quando ele cai estritamente
dentro desse intervalo e usa bissecção nos outros casos; a convergência é testada ângulo a ângulo,
e, se o limite de iterações for atingido, o cálculo termina com bissecção pura. Comparada à
bissecção em 2.030 geometrias, nos dois planos e de 0 a $85^\circ$, a versão final deu diferença
máxima de $1{,}2\times10^{-8}$~m em GZ, e as curvas ficaram sem picos.

\subsection{Listagens do código}

As listagens a seguir são cópias exatas do arquivo \path{caso_base/codigo/pnv3425_grupo02.py}, com a
numeração de linhas original, extraídas por \path{relatorio/gerar_trechos.py}. O arquivo da condição
de 10\% de lastro é idêntico, exceto pelas linhas da \lst{diff}.

\listagem{dados}{Dados fixos, coeficientes do Código MODU, limites das restrições e hipóteses parametrizadas}
\listagem{cenario}{Cenário de projeto, plataforma de referência e domínio das variáveis}
\listagem{hidrostatica}{Hidrostática (Exercício 7)}
\listagem{pesos}{Pesos, centros e altura metacêntrica (Exercício 8)}
\listagem{caixas}{Retângulos da projeção 2D e integrais exatas da parte submersa}
\listagem{equilibrio}{Linha d'água a volume constante por Newton com salvaguarda e curva GZ}
\listagem{alagamento}{Ângulo de alagamento}
\listagem{vento}{Centro de resistência lateral e tabela de vento por plano}
\listagem{criterio}{Critério de áreas e laço de estabilidade nos dois planos}
\listagem{heave}{Massa adicional, espectro JONSWAP e resposta em \emph{heave} (Exercício 10)}
\listagem{avaliar}{Modelo completo, \emph{air gap} mínimo e tabela de restrições}
\listagem{lhs}{Filtro de não dominância e variação sistemática por hipercubo latino}
\listagem{problema}{\emph{Cache} do problema e partida do SLSQP}
\listagem{somapond}{Soma ponderada com \emph{multi-start} e continuação}
\listagem{epsilon}{Método $\varepsilon$-restrito e escolha do joelho}
\listagem{fronteira}{Montagem da fronteira de Pareto com reconfirmação dos candidatos do LHS}

\subsection{Diferenças da condição de 10\% de lastro}

A \lst{diff}, mostra o \emph{diff} unificado entre o código do caso-base e o da condição de 10\% de
lastro. Linhas iniciadas por ``$-$'' existem só no caso-base, linhas iniciadas por ``+'' só na
condição de 10\%, e os cabeçalhos \texttt{@@} indicam as linhas de cada arquivo. A única mudança de
cálculo é o valor padrão de \texttt{LASTRO\_MIN\_FRAC}, que ativa a restrição
$V_{lastro} \ge 0{,}10\,V_{pontoons}$ em \texttt{tabela\_restricoes()}; as demais trocam a pasta de
saída e os textos.

\par\needspace{10\baselineskip}\medskip
\captionof{codigo}[Diferenças entre o código do caso-base e o da condição de 10\% de lastro]{Diferenças
entre \texttt{caso\_base/codigo/pnv3425\_grupo02.py} e
\texttt{condicao\_10pct\_lastro/codigo/pnv3425\_grupo02\_condicao\_10pct\_lastro.py}.}\label{lst:diff}
\VerbatimInput[numbers=none]{trechos/diff_lastro10.txt}

\subsection{Como reproduzir os resultados}

Os comandos abaixo refazem todos os resultados e o relatório, a partir da raiz do repositório.

\begin{Verbatim}[numbers=none]
pip install -r requirements.txt

cd caso_base
python codigo/pnv3425_grupo02.py                  # ~4 min em 8 núcleos
python codigo/analises_complementares.py
python codigo/comparativo_plataformas_reais.py
python codigo/comparativo_poster_2011.py
python codigo/figuras_apresentacao.py

cd ../condicao_10pct_lastro
python codigo/pnv3425_grupo02_condicao_10pct_lastro.py   # ~3 min em 8 núcleos
python codigo/analises_condicao_10pct_lastro.py           # usa ../caso_base/resultados

cd ../relatorio
python gerar_trechos.py                           # trechos de código do Apêndice A
tectonic relatorio.tex                            # ou XeLaTeX
\end{Verbatim}

As hipóteses alternativas da Seção~\ref{sec:sensibilidade} são selecionadas por variáveis de
ambiente, com \texttt{OUT\_DIR} apontando para outra pasta: \texttt{HM=15}, \texttt{INCLUIR\_GAS=1}
ou \texttt{LASTRO\_MIN\_FRAC=0.10}. Os modelos aceitam também \texttt{HD}, \texttt{PASSO\_W},
\texttt{N\_AMOSTRAS}, \texttt{N\_EPSILON}, \texttt{N\_PROCESSOS}, \texttt{REAPROVEITAR=1}, que
retoma etapas já salvas, e \texttt{MOSTRAR\_GRAFICOS=0}. O relatório usa \texttt{fontspec} e deve
ser compilado com XeLaTeX ou Tectonic; no Overleaf, é preciso enviar o repositório inteiro, porque
as figuras são lidas das pastas de resultados.

% =============================================================================
\section{Fronteiras de Pareto completas}

Este apêndice lista todos os pontos não dominados das duas versões do projeto, ordenados por peso.
Na última coluna, SP indica a soma ponderada e $\varepsilon$ o método $\varepsilon$-restrito.

\subsection{Caso-base}\label{ap:fronteira}

\begin{longtable}{rrrrrrrrl}
\caption{Pontos não dominados do caso-base.}\label{tab:fronteira}\\
\toprule
\# & $L_p$ & $b_p$ & $h_p$ & $b_c$ & $h_c$ & $T$ & $\Wc$ [t] & $\Hz$ [m] / origem \\
\midrule
\endfirsthead
\toprule
\# & $L_p$ & $b_p$ & $h_p$ & $b_c$ & $h_c$ & $T$ & $\Wc$ [t] & $\Hz$ [m] / origem \\
\midrule
\endhead
1 & 137,49 & 16,39 & 7,37 & 8,00 & 20,00 & 12,00 & 10.346 & 2,366 / SP (cen. 1) \\
2 & 135,76 & 13,06 & 9,38 & 8,00 & 20,00 & 13,89 & 10.358 & 2,265 / $\varepsilon$ \\
3 & 135,44 & 12,18 & 10,08 & 8,00 & 20,00 & 14,54 & 10.362 & 2,230 / SP \\
4 & 135,49 & 9,33 & 13,40 & 8,00 & 20,00 & 16,12 & 10.530 & 2,163 / $\varepsilon$ \\
5 & 136,11 & 8,00 & 15,71 & 8,00 & 20,00 & 17,48 & 10.622 & 2,105 / SP \\
6 & 139,19 & 8,00 & 15,42 & 8,00 & 22,07 & 17,45 & 10.796 & 2,061 / $\varepsilon$ \\
7 & 139,22 & 8,00 & 15,41 & 8,00 & 22,09 & 17,45 & 10.798 & 2,061 / SP \\
8 & 145,89 & 8,00 & 14,83 & 8,00 & 26,65 & 17,37 & 11.190 & 1,979 / SP \\
9 & 147,72 & 8,00 & 14,68 & 8,00 & 27,93 & 17,34 & 11.301 & 1,959 / $\varepsilon$ \\
10 & 149,67 & 8,00 & 14,53 & 8,00 & 29,30 & 17,33 & 11.420 & 1,939 / SP \\
11 & 150,00 & 13,22 & 8,50 & 8,76 & 37,84 & 17,15 & 12.239 & 1,858 / $\varepsilon$ \\
12 & 150,00 & 13,58 & 8,31 & 8,96 & 39,91 & 17,30 & 12.606 & 1,846 / SP \\
13 & 150,00 & 13,91 & 8,16 & 9,50 & 45,85 & 18,75 & 13.657 & 1,756 / $\varepsilon$ \\
14 & 108,06 & 13,31 & 10,70 & 12,49 & 35,15 & 23,00 & 14.241 & 1,654 / $\varepsilon$ \\
15 & 111,80 & 13,40 & 10,21 & 12,55 & 39,61 & 23,98 & 14.998 & 1,552 / $\varepsilon$ (cen. 3) \\
16 & 116,45 & 13,55 & 9,67 & 12,60 & 44,92 & 25,02 & 15.934 & 1,451 / $\varepsilon$ \\
17 & 117,68 & 14,06 & 9,14 & 13,10 & 50,00 & 26,01 & 17.436 & 1,369 / SP \\
18 & 117,13 & 14,01 & 9,74 & 13,22 & 50,00 & 26,90 & 18.067 & 1,349 / $\varepsilon$ \\
19 & 115,14 & 13,62 & 12,56 & 13,61 & 50,00 & 31,26 & 20.641 & 1,247 / $\varepsilon$ \\
20 & 115,39 & 13,60 & 12,89 & 13,60 & 50,00 & 31,63 & 20.912 & 1,239 / SP \\
21 & 107,53 & 16,62 & 10,64 & 14,70 & 50,00 & 36,16 & 21.934 & 1,145 / $\varepsilon$ \\
22 & 107,53 & 15,78 & 12,92 & 14,59 & 50,00 & 40,00 & 23.327 & 1,046 / SP \\
23 & 107,53 & 15,85 & 12,87 & 14,59 & 50,00 & 40,00 & 23.346 & 1,045 / SP \\
24 & 109,17 & 14,77 & 14,25 & 14,35 & 50,00 & 40,00 & 23.531 & 1,044 / SP \\
25 & 110,11 & 14,22 & 15,03 & 14,22 & 50,00 & 40,00 & 23.634 & 1,043 / SP (cen. 2) \\
\bottomrule
\end{longtable}

\clearpage
\subsection{Condição de 10\% de lastro}\label{ap:fronteiraL}

\begin{longtable}{rrrrrrrrl}
\caption{Pontos não dominados na condição de 10\% de lastro (gerada).}\label{tab:fronteiraL}\\
\toprule
\# & $L_p$ & $b_p$ & $h_p$ & $b_c$ & $h_c$ & $T$ & $\Wc$ [t] & $\Hz$ [m] / origem \\
\midrule
\endfirsthead
\toprule
\# & $L_p$ & $b_p$ & $h_p$ & $b_c$ & $h_c$ & $T$ & $\Wc$ [t] & $\Hz$ [m] / origem \\
\midrule
\endhead
\linhas{tabelas_lastro10/tab_fronteira_linhas.tex}
\bottomrule
\end{longtable}

\end{document}
